\documentclass[11pt,a4paper,reqno]{amsart}
\usepackage[T1]{fontenc}
\usepackage{lmodern,mathtools,amssymb,amsthm,mathrsfs,booktabs}
\usepackage[expansion=false]{microtype}
\usepackage[a4paper,textwidth=156mm,textheight=237mm,centering]{geometry}
\usepackage[hidelinks]{hyperref}
\hypersetup{pdftitle={Fourier decay for shrinking-target centers of hyperbolic toral automorphisms},pdfauthor={Nikita Maximilian Grimm}}
\pdfinfoomitdate=1
\pdftrailerid{}
\pdfsuppressptexinfo=15
\newtheorem{theorem}{Theorem}[section]
\newtheorem{lemma}[theorem]{Lemma}
\newtheorem{proposition}[theorem]{Proposition}
\newtheorem{corollary}[theorem]{Corollary}
\theoremstyle{definition}
\newtheorem{definition}[theorem]{Definition}
\newtheorem{example}[theorem]{Example}
\theoremstyle{remark}
\newtheorem{remark}[theorem]{Remark}
\numberwithin{equation}{section}
\setlength{\emergencystretch}{1.5em}
\raggedbottom
\title[Fourier decay for shrinking-target centers]{Fourier decay for shrinking-target centers of hyperbolic toral automorphisms}
\thanks{Curated by Nikita Maximilian Grimm with AI assistance. Repository: \url{https://github.com/NikitaMGrimm/math}.}
\date{}
\begin{document}
\begin{abstract}
For a hyperbolic automorphism of the two-torus, we study the centers of
exponentially shrinking balls hit infinitely often from a dense set of
parameters on a prescribed arc. The unstable eigenline, a mixed line and
a curved arc give three different sharp Fourier decay scales at the same
Hausdorff dimension. A straight subarc can distinguish centers with dense
fibers from those with nonempty fibers. We also obtain sparse-time criteria
and a critical center whose fibers on every transverse line exclude a
prescribed decay rate.
\end{abstract}
\maketitle

\section{Centers and hitting parameters}\label{sec:intro}
Let \(A\in{\rm GL}_2(\mathbb Z)\) be hyperbolic, let \(\lambda>1\) be its
spectral radius, and fix \(c>0\). For a parametrized arc
\(\gamma:I\to\mathbb T^2=\mathbb R^2/\mathbb Z^2\) on a nondegenerate
compact interval and times \(S\subset\mathbb N\), put
\[
 E_y=\{x\in I:\operatorname{dist}(A^n\gamma(x),y)<r_n
                  \text{ for infinitely many } n\in S\},
 \qquad r_n=c\lambda^{-\tau n}.
\]
Denote by \(D\) the centers with \(E_y\) dense in \(I\), by \(P\) those with
\(E_y\ne\varnothing\), and by \(V\) the centers within \(r_n\) of
\(A^n\gamma(I)\) infinitely often. Thus \(D\subseteq P\subseteq V\).
Only \(V\) allows the witness to change at every visit.

\begin{theorem}[Three sharp scales]\label{thm:planar}
Let \(\tau>1\), \(S=\mathbb N\), and \(d_0=1+1/\tau\).
For each geometry below, \(D,P,V\) have Hausdorff dimension \(d_0\) in
every nonempty open target patch and Fourier dimension zero. Write \(L=L(\xi)=
\log(e+|\xi|)\). In every patch there is a probability on a compact
subset of \(D\) attaining the middle column at all real frequencies in a
local lift, with finite \(s\)-energy simultaneously for all \(0<s<d_0\).
For any \(\epsilon>0\), every finite positive torus measure satisfying the
last column on integer frequencies gives \(V\) zero mass.
\[
\begin{array}{lll}
\toprule
\text{Initial arc} & \text{Attainable decay} & \text{Zero-mass bound}\\
\midrule
\text{Affine unstable line} & o((\log L)^{-1}) & O((\log L)^{-1-\epsilon})\\
\text{Affine mixed line} & o(L^{-1}) & O(L^{-1-\epsilon})\\
\text{Curved } C^2\text{ arc} & o(L^{-3/2}) & O(L^{-3/2-\epsilon})\\
\bottomrule
\end{array}
\]
A mixed direction has two nonzero eigencomponents. In the third row the
arc is regular and its curvature is nonzero somewhere in every open
parameter interval. Isolated tangencies to either eigendirection are allowed.
\end{theorem}

The dimensions in Theorem~\ref{thm:planar} persist for every infinite
\(S\), even though each individual fiber \(E_y\) has dimension zero
(with \(\dim_H\varnothing=0\)). With
\(U_n(J)=\bigcup_{x\in J}B(A^n\gamma(x),r_n)\) and \(\mathcal J\)
a countable basis of open intervals in the interior of \(I\), one has
\begin{equation}\label{eq:dense}
 D=\bigcap_{J\in\mathcal J}\bigcap_{N\ge1}
       \bigcup_{\substack{n\in S\\ n\ge N}}U_n(J).
\end{equation}
Indeed every open tail union in parameter space is dense for such \(y\).
Their intersection is dense by Baire's theorem. This identity allows the
construction to visit a different parameter interval at each stage while
retaining infinitely many visits by each parameter in a dense set.

A \emph{regular gauge} \(h\) is positive and bounded on \([1,\infty)\),
nonincreasing and tending to zero, with \(h(2t)\asymp h(t)\) at infinity;
eventual monotonicity suffices after a bounded initial redefinition. Write
\(H(\xi)=h(L(\xi))\), and let \(\mathcal B(S)\) be the indices of the
dyadic time blocks \([2^j,2^{j+1})\) meeting \(S\). Fourier transforms use
\(e(t)=\exp(2\pi i t)\) and the negative sign in the exponent. We use
\(I_s(\mu)=\iint |y-z|^{-s}\,d\mu(y)d\mu(z)\) in a local Euclidean chart. Fourier
dimension is the supremum of \(s\) up to the ambient dimension for which a
nonzero measure carried by the set has decay \(O((1+|\xi|)^{-s/2})\).

\begin{theorem}[Available times]\label{thm:sparse}
Assume \(\tau>1\), use the three arc classes of Theorem~\ref{thm:planar},
and set \(d_0=1+1/\tau\). Fix any infinite \(S\) and a regular \(h\).
For the unstable and mixed lines the respective tests are
\begin{equation}\label{eq:tests}
        \sum_{j\in\mathcal B(S)}h(2^j),\qquad \sum_{n\in S}h(n).
\end{equation}
Divergence gives a compactly supported probability as in
Theorem~\ref{thm:planar} with decay \(o(H)\) and the same energies.
Convergence forces every finite positive measure with decay \(O(H)\)
to give \(V\) zero mass.

For curved arcs, convergence of \(\sum_{n\in S}h(n)^{2/3}\) forces the
same zero-mass conclusion. Either of the following conditions gives the
compact probability with decay \(o(H)\) and all the stated energies:
\begin{enumerate}
\item \(\sum_{n\in S}h(n)=\infty\);
\item for each \(K>0\) there is \(T\subseteq S\) with
\(\sum_{n\in T}h(n)^{2/3}=\infty\) and \(m-n\ge K\log(e+n)\)
at successive \(n<m\) in \(T\).
\end{enumerate}
In particular, \(\sum_{n\in S}h(n)^{2/3}/\log(e+n)=\infty\)
suffices. All three local Hausdorff dimensions are \(d_0\).
\end{theorem}

For all times the first test in \eqref{eq:tests} is equivalent to
divergence of \(\sum_n h(n)/n\) by dyadic summation. For the curved
endpoint \(h(n)=n^{-3/2}\), the last sufficient sum is the divergent
harmonic logarithmic series. The distinction between \(D\) and \(P\)
becomes visible on a single regular \(C^2\) arc: in eigenvector coordinates
take a small positive dilation of \(t\mapsto(t,t+\varepsilon(t-1/2)_+^3)\)
on \([0,1]\) with \(\varepsilon>0\). Its straight half forces the sharp logarithmic power
on \(D\) to be 1. Its curved half attains power \(3/2\) on \(P\) and \(V\).
We prove these assertions in Section~\ref{sec:curves}.

\begin{theorem}[One critical center]\label{thm:critical}
Fix \(A\), an infinite \(S\), and a positive nonincreasing function
\(g(t)\to0\). There exists a center \(y\) such that for every affine
segment transverse to the stable direction the parameter fiber along \(S\)
at radius \(c\lambda^{-n}\), for every \(c>0\), is a dense \(G_\delta\)
of dimension zero. Every nonempty open subinterval of the parameter interval
contains a compact subset of the fiber carrying a Rajchman probability.
Every finite positive measure on the parameter line with
\(\widehat\mu(\xi)=O(g(L(\xi)))\) gives the fiber zero mass. The same
\(y\) works for open eigenrectangles of stable half-width \(c_s\lambda^{-n}\)
and unstable half-width \(c_u e^{-\alpha n}\), for all \(c_s,c_u,\alpha>0\).
\end{theorem}

A Rajchman measure has Fourier transform tending to zero. The measure
in Theorem~\ref{thm:critical} lives in parameter space. The first two
theorems construct measures in target space. The change of variable is
substantial: the centers hit from a single fixed parameter have dimension
zero for every positive exponential shrinking rate, simply by covering
with one ball per level.

\subsection{Relation to earlier work}
Hill and Velani~\cite{HV} introduced matrix shrinking-target dimension
problems on tori. Hu, Persson, Wu and Zhang~\cite[Theorem~2.1 and
Remark~2.2]{HPWZ} give the planar saddle law for ambient starting points
and describe critical dependence on moving centers in the unimodular case.
Their local subcritical ball result gives the transverse slice dimension
by stable saturation, as explained below. General mass transference
principles concern Hausdorff measure of neighborhoods of resonant sets
under local scaling assumptions~\cite{BV,AB}.

There is a direct target-space antecedent when the line through the origin
traverses a whole closed rational geodesic. At level \(n\) its image is the
kernel modulo one of a primitive integer normal \(q_n\), and the swept
event is \(\|q_n\cdot y\|<|q_n|r_n\), where \(\|\cdot\|\) denotes distance
to the nearest integer. In the unimodular case \(|q_n|\asymp\lambda^n\).
Squeezing the right side between powers of \(|q_n|\) with exponents
\(1-\tau\pm\epsilon\), Rynne's restricted-vector theorem
\cite[p.~70]{Rynne} gives \(\dim_H V=1+1/\tau\). The upper-bound
argument of Cai and Hambrook~\cite[Theorem~1.4.1 and \S3.3]{CH}
applies to measures carried by \(V\) and gives \(\dim_F V=0\):
\(\sum_n |q_n|r_n<\infty\) and \(\sum_n r_n^s<\infty\) for every \(s>0\).
These dimension calculations for \(V\) leave open the rates of decay and
the effect of requiring a dense parameter fiber.

Chow and Zhou~\cite[Theorems~4, 8, 11 and Corollary~7]{CZ} vary the
target for a fixed starting point and prove metric laws for almost every
starting point. Their nonsingular-differences condition in Theorems~8
and~11 holds here: \(A^m(A^{n-m}-I)\) is invertible for \(n>m\).
Their Theorem~4 uses growth of the least singular values of differences,
and Corollary~7 uses successive expansion and Fourier decay of the
starting measure. A stable eigenvector prevents this expansion for saddle
powers. Our positive estimates concern a measure on centers for a
prescribed initial arc.

For the Fourier obstructions, Pollington, Velani, Zafeiropoulos and
Zorin~\cite[Theorem~2]{PVZZ} give a scalar fixed-shift convergence test
for summable widths and summable Fourier suprema. Salem's covering
criterion, as presented by Kahane~\cite[Ch.~VII, \S8, p.~97]{Kahane},
uses simultaneous approximation of interval centers. Lyons
\cite[Theorems~4--5, p.~608, and Lemma~9, pp.~613--614]{Lyons}
proved the inverse-log-log obstruction for lacunary sequences and
attained this scale on the set of base-two non-normal numbers.
Smooth product constructions of measures on approximation sets
go back to Kaufman and Bluhm~\cite{Kaufman,Bluhm}; see also
Fraser and Nguyen~\cite{FN} and Polasek and Rela~\cite{PR}.
The finite phase majorant and weighted averaging argument used for power
targets in~\cite{Power} will be proved here in the forms needed for lines.
For curved arcs a positive correction to those averages is needed to
reach the exponent \(3/2\).

\section{Measures on centers with dense fibers}\label{sec:construction}
We work on \(\mathbb T^m\), with normalized Haar measure \(dy\). It
suffices to construct measures in arbitrarily small patches. In a chart
of diameter less than \(1/4\) Euclidean and torus distances coincide.
For smooth densities \(F\), use the two norms
\[
 \|F\|_H=\sup_{k\in\mathbb Z^m}\frac{|\widehat F(k)|}{H(k)},\qquad
 \|F\|_s^2=\sum_{k\in\mathbb Z^m}(1+|k|)^{s-m}|\widehat F(k)|^2,
 \qquad 0<s<m.
\]
The latter squared norm is comparable, for positive measures, to energy
with kernel \(1+\operatorname{dist}(x,y)^{-s}\). Indeed integrate the
periodized heat kernel at times \(0<t<1\) against \(t^{(m-s)/2-1}dt\).
Termwise Fourier transformation gives coefficients comparable to the
weights above. In its Gaussian representation, substitution
\(u=\operatorname{dist}(x,y)^2/t\) in the nearest translate gives the
singularity of order \(s\); the remaining translates give a bounded
contribution away from lattice points. Truncation of the time integral
and monotone convergence justify the comparison for measures.

\begin{lemma}[Localization and compact support]\label{lem:compact}
Suppose \(0<d_0\le m\). For every basis interval \(J\), time tail,
\(s\in(0,d_0)\) and \(\delta>0\), suppose there is a smooth nonnegative
\(Q\) supported in a finite union of \(U_n(J)\) from that tail with
\(\|Q-1\|_H+\|Q-1\|_s<\delta\). In every open target patch there is a
probability on a compact subset of \(D\) with \(\widehat\mu(\xi)=o(H(\xi))\)
and finite \(I_s\) simultaneously for \(0<s<d_0\).
\end{lemma}
\begin{proof}
Multiplication by a fixed smooth function \(F\) is bounded on each norm.
For the first norm use \(H(k-\ell)/H(k)\ll(1+|\ell|)^C\), from doubling
and \(L(k)\le L(k-\ell)+\log(1+|\ell|)\). For the second norm the ratio
of the square-root weights has polynomial growth in \(\ell\), so Young's
inequality applies to convolution with the rapidly decreasing coefficients
of \(F\).

Start with a smooth probability in a compact subpatch. Let \(s_j\uparrow d_0\)
and choose a sequence of basis intervals containing every member
infinitely often. At step \(j\) replace the previous density \(F\) by
\(FQ/\int FQ\), using the next interval and levels later than all earlier
ones. Evaluation at zero bounds the change in mass. By choosing \(\delta\)
after \(F,s_j\), the mass stays positive and normalization leaves errors
at most \(2^{-j}\) in both norms. The new smooth densities have nested
compact supports. The transforms converge in the weighted supremum norm,
so their limit divided by \(H\) vanishes at infinity. A weak cluster
probability has exactly these coefficients. The \(s_j\)-norm bounds the
\(s\)-norm whenever \(s\le s_j\), hence the summable errors and Fatou give
every asserted energy. The limit support meets the condition in
\eqref{eq:dense} point by point, since it lies in every stage support.

For \(\xi=k+t\) with \(t\in[0,1]^m\), multiply the measure in its local
lift by a chart cutoff equal to one on the support times \(e(-t\cdot y)\).
Its multiplier has uniformly rapidly decreasing Fourier coefficients.
Split convolution into a bounded range of shifts and its uniformly
summable tail to transfer the little-\(o\) bound to real \(\xi\).
\end{proof}

\begin{lemma}[Averaging]\label{lem:average}
Fix \(J\). Let \(G_n\) be smooth probability densities supported in
\(U_n(J)\). Suppose their ball masses satisfy
\(\int_{B(y,R)}G_n\ll_J R^{d_0}\) for \(0<R\le1\). If bounded
nonnegative weights of divergent sum satisfy
\begin{equation}\label{eq:weighted}
 \sum_{n\in S}w_n\big|\widehat G_n(k)-{\bf1}_{k=0}\big|
       \ll_J H(k)\qquad(k\in\mathbb Z^m),
\end{equation}
then these levels provide the multipliers in Lemma~\ref{lem:compact}.
\end{lemma}
\begin{proof}
A normalized finite average \(Q=\sum w_nG_n/\sum w_n\) in a fresh tail
has \(\|Q-1\|_H\ll_J1/\sum w_n\). For \(s<s'<d_0\) its ball bound gives
a uniform \(s'\)-energy bound by integration in the radius. Above
frequency \(K\) the squared \(s\)-norm of the error is \(O_J(K^{s-s'})\).
Take \(K\) large, then the sum of weights large to control the remaining
finite set of modes.
\end{proof}

\subsection{Planar tube measures}
Use unit right eigenvectors \(e_u,e_s\) for the real signed eigenvalues
\(a,b\), where \(|a|=\lambda\), \(|b|=\lambda^{-1}\).
A complex conjugate pair would have modulus one. A rational
eigenvalue would be an integral unit, so the characteristic polynomial
is irreducible. Write \(z=z_u e_u+z_s e_s\) in primal coordinates and
write \(k_u=k\cdot e_u,\ k_s=k\cdot e_s\) for integer characters.
The two coordinate systems are dual even when the eigenvectors are oblique.
For nonzero \(k,z\in\mathbb Z^2\),
\begin{equation}\label{eq:norm}
             |k_u k_s|\gg_A1,\qquad |z_u z_s|\gg_A1.
\end{equation}
Indeed \(\det[z,Az]\) is a fixed nonzero multiple of \(z_u z_s\),
and is a nonzero integer by irreducibility. Apply the same argument to
\(A^{\mathsf T}\) in the dual eigenbasis for the character estimate.

On a graph portion \(u\mapsto u e_u+f(u)e_s\) of an arc, choose a
probability \(\sigma\) with a smooth compactly supported density in \(u\).
For a nonnegative smooth probability \(\chi\) supported in the Euclidean
ball of radius \(1/4\), write \(\chi_r\) for its periodized rescaling.
Define
\begin{equation}\label{eq:tube}
 G_n=\chi_{r_n}*(A^n_*\sigma),\qquad
 \widehat G_n(k)=\widehat\chi(r_n k)\widehat\sigma((A^n)^{\mathsf T} k).
\end{equation}
Here the transform of \(\chi\) on the right is Euclidean. The support
lies compactly inside the tube from the chosen parameter interval.

\begin{lemma}[Ball mass of a graph tube]\label{lem:ball}
For \(\tau\ge1\) these densities have ball mass \(O_\sigma(R^{1+1/\tau})\)
for every radius \(0<R\le1\), uniformly in \(n\) and the center.
\end{lemma}
\begin{proof}
The possible integer translates of a ball of radius \(R\) met by the
unsmoothed graph lie, in primal coordinates, in a translated box with
extents \(O(\lambda^n)\) and \(O(R+\lambda^{-n})\). Add to each lattice
point a translate of the unit-area fundamental parallelogram
\(A^n[0,1)^2\), whose coordinate extents are \(O(\lambda^n)\) and
\(O(\lambda^{-n})\). The interiors are disjoint, so area bounds the count
by \(O(1+R\lambda^n)\). Each passage has parameter mass
\(O_\sigma(R\lambda^{-n})\) by projection on the unstable coordinate.
Thus the total mass is \(O_\sigma(R^2+R\lambda^{-n})\).
For \(R\ge r_n\) smoothing costs only a fixed enlargement, and
\(\lambda^{-n}\ll R^{1/\tau}\). At smaller \(R\) the unsmoothed estimate
at \(r_n\) bounds the smoothed density by \(O_\sigma(r_n^{1+1/\tau-2})\);
multiplication by \(O(R^2)\) finishes the proof.
\end{proof}

\subsection{Dimension and localization in the parameter}
The image of any fixed compact \(C^1\) arc has length \(O(\lambda^n)\).
Its tube is covered by \(O(1+\lambda^n/r_n)\) balls of radius \(2r_n\).
The resulting Hausdorff sum gives the upper dimension \(d_0\) for \(V\).
Conversely any infinite schedule admits a regular \(h\) for which even
the occupied-block sum in \eqref{eq:tests} diverges: take successive
ranges of dyadic indices containing at least \(2^i\) occupied blocks
and set \(h=2^{-i}\) on the corresponding time ranges. Doubling crosses
at most one boundary. The positive constructions below with this \(h\)
give local dimension at least \(d_0\) on \(D\) by finite energy. A convergent
gauge such as \(h(t)=(1+t)^{-3}\) bounds every positive power of Fourier
decay, so the converses give Fourier dimension zero. For a fixed center
when \(\tau>1\), pulling back its ball gives width
\(O(\lambda^{-(\tau-1)n})\) even along the longer inverse axis.
Only boundedly many lifts can meet the bounded lifted arc. A regular
compact \(C^1\) parametrization has a finite cover by locally bi-Lipschitz
charts. The exponential diameter bound therefore gives dimension zero
for each individual parameter fiber.

\section{Affine lines and integer normals}\label{sec:lines}
Write \(\gamma(x)=p+xv\) with \(v_u\ne0\). A single normal direction
can test many consecutive tubes when \(v_s=0\). For affine translates
the scalar phases can still change at every time.

\begin{lemma}[A finite phase majorant]\label{lem:majorant}
Let \(z_1,\ldots,z_m\in\mathbb T\), \(0<\eta<1/16\), and
\(1\le M\le(16\eta)^{-1}\) an integer. For an integer character
\(k\ne0\) of a torus and a finite positive measure \(\nu\) on it,
\[
 \nu\{y:\min_i\|k\cdot y-z_i\|<\eta\}
 \ \ll\ 16^m M^{-1}\|\nu\|
             +\sup_{1\le j<M}|\widehat\nu(jk)|.
\]
The estimate holds also with \(k\cdot y\) replaced by \(Qx\) on
the real line, using frequencies \(jQ\). For \(M=1\) take the supremum
to be zero.
\end{lemma}
\begin{proof}
Partition \(\mathbb T^m\) into \(16^m\) cubes of side \(1/16\).
At least \(M/16^m\) of the indices \(j=1,\ldots,M\) put
\(j(z_1,\ldots,z_m)\) in the same cube. Denote them by \(F\). For
\(t\) in any test interval the phases \((j-j_0)t\), \(j,j_0\in F\),
are within \(1/8\) of zero modulo one. Thus
\(\left||F|^{-1}\sum_{j\in F}e(jt)\right|^2\ge1/2\).
This polynomial has nonnegative coefficients of sum one and constant
term \(1/|F|\). Integration at \(t=k\cdot y\), or at \(t=Qx\), gives
the estimate.
\end{proof}

\subsection{The obstruction}
For a mixed line take \(0<\delta<\tau-1\). For any \(w\ne0\) and \(T>2\),
Dirichlet approximation to the ratio of its smaller coordinate to its
larger gives \(k\in\mathbb Z^2\setminus\{0\}\) with \(|k|\ll T\) and
\(|k\cdot w|\ll |w|/T\). Apply this with \(w=A^nv\) and
\(T=\lambda^{(1+\delta)n}\). For \(|k|\le\lambda^{n/3}\) the fast
term \(|a^nv_uk_u|\gg_v\lambda^{2n/3}\) dominates the slow term
by \eqref{eq:norm}; hence our normal has \(|k|>\lambda^{n/3}\) eventually.
Throughout \(U_n(I)\),
\[
 \|k\cdot(y-A^np)\|\ll_I
          \lambda^{-\delta n}+\lambda^{-(\tau-1-\delta)n}.
\]
Lemma~\ref{lem:majorant} with one phase and \(M\) comparable to the
inverse width bounds the mass by \(O(e^{-cn}+h(n))\) under the stated
Fourier hypothesis (here and below exponential savings may use different
positive constants). Summing over \(S\) proves the mixed converse.

When \(v=v_u e_u\), fix \(\kappa=(\tau-1)/2\). Apply the same
Dirichlet approximation to \(e_u\) at \(T=\lambda^{(1+\kappa)N}\).
The resulting \(k\) has \(|k|\asymp T\) by \eqref{eq:norm}.
For \(N\le n\le(1+\delta)N\), its error on \(U_n(I)\) is
\[
 \|k\cdot(y-A^np)\|\ll_I
      \lambda^{-(\kappa-\delta)N}+\lambda^{-\kappa N}.
\]
Choose \(\delta<\kappa/2\) small enough that
\(\delta\log16<\kappa\log\lambda/4\). Now one character tests the
whole block with at most \(\lfloor\delta N\rfloor+1\) phases.
Lemma~\ref{lem:majorant} with \(M=\lfloor\lambda^{\kappa N/3}\rfloor\)
bounds the block mass by \(O(e^{-c'N}+h(N))\). Each occupied dyadic
block costs only boundedly many such shorter blocks with comparable
starts. This proves the occupied-block converse by Borel--Cantelli.

\subsection{Supporting measures}
Take a smooth parameter density \(f_J\) of integral one compactly
supported in \(J\). The coefficients of the tube density \eqref{eq:tube}
are
\begin{equation}\label{eq:linefourier}
 e(-k\cdot A^np)\,\widehat\chi(r_nk)\,
 \widehat f_J(a^n v_u k_u+b^n v_s k_s).
\end{equation}
For \(0<|k|\le\lambda^{n/3}\) the fast term gives an exponentially
growing gap as above. For \(|k|\ge \lambda^{(\tau+1)n}\) the bump gives
\(O((1+|k|)^{-3})\) by taking sufficiently many derivatives. The two
outer ranges therefore contribute \(O(H(k))\) to the sum of absolute
coefficients with any bounded weights: in the first range the exponential
tail starts at \(n\gg L(k)\); in the second there are \(O(L(k))\)
terms. Regularity gives \(h(t)\gg t^{-C}\).

In the intervening range \(n\asymp L(k)\). On an unstable line use
one representative of each occupied dyadic block, weighted by \(h(n)\).
There are \(O(1)\) representatives in this middle range. On a mixed line
use all available levels with weight \(h(n)\). To bound their contribution
put \(X=v_uk_u,\ Y=v_sk_s\), and \(t_*=\tfrac12\log_\lambda|Y/X|\).
The larger of \(|a^n X|\) and \(|b^nY|\) equals
\(\sqrt{|XY|}\lambda^{|n-t_*|}\), with \(|XY|\gg_v1\).
For \(|n-t_*|\ge1\) their signed sum has modulus at least
\(1-\lambda^{-2}\) times this maximum. The bounded number of remaining
terms and rapid decay away from this balance time give
\(\sup_{k\ne0}\sum_n|\widehat f_J(a^nX+b^nY)|<\infty\).
Since \(h(n)\asymp H(k)\) on the middle range, \eqref{eq:weighted}
holds in both cases. Lemmas~\ref{lem:ball} and~\ref{lem:average}
complete the line assertions of Theorems~\ref{thm:planar}
and~\ref{thm:sparse}.

\section{Curved arcs}\label{sec:curves}
Keep the planar eigenvector coordinates of Section~\ref{sec:construction}.
In particular, the dual coordinates of \((A^n)^{\mathsf T}k\) are
\((a^n k_u,b^n k_s)\), and \eqref{eq:norm} gives
\(\min(|k_u|,|k_s|)\gg_A |k|^{-1}\) for each nonzero integer \(k\).
The eigenvalues may have either sign. Fix any \(c>0\), and write
\(d_0=1+1/\tau\). We use \(\tau>1\) for the obstruction and
\(\tau\ge1\) for the constructions.

We will use two consequences of gauge regularity. Iterating
\(h(2t)\ge c_h h(t)\) and using monotonicity between dyadic arguments
gives \(h(t)\gg (1+t)^{-\beta}\) for some \(\beta>0\) and all large
\(t\). Also, \(h(t)\asymp h(t')\) when \(t,t'\) are large and their
ratio stays between two positive constants. Both comparisons concern
only the tail of \(h\).

\subsection{A majorant for the entire arc}
\begin{lemma}[Fourier mass of a smoothed neighborhood]\label{lem:arc-majorant}
Let \(\Gamma\subset\mathbb T^2\) be the image of a rectifiable arc.
For all sufficiently small \(\rho>0\), there is a smooth nonnegative
\(F_\rho\) that equals one on the open \(\rho\)-neighborhood of \(\Gamma\)
and satisfies
\begin{equation}\label{eq:curve-majorant}
 \widehat F_\rho(0)\ll \rho,\qquad
 \sum_{\ell\in\mathbb Z^2}|\widehat F_\rho(\ell)|\ll \rho^{-1/2},
 \qquad
 \sum_{|\ell|>X}|\widehat F_\rho(\ell)|
       \ll_M \rho^{-1/2}(1+\rho X)^{-M}
 \quad(X\ge0).
\end{equation}
\end{lemma}
\begin{proof}
Write \(N_t(\Gamma)\) for the open \(t\)-neighborhood. An arclength cover
uses \(O_\Gamma(\rho^{-1})\) balls of radius \(O(\rho)\), so
\(|N_{2\rho}(\Gamma)|\ll_\Gamma\rho\). For a smooth probability bump
\(\psi\) supported in the Euclidean unit ball, let
\(F_\rho=\mathbf 1_{N_{2\rho}(\Gamma)}*\psi_{\rho/2}\), with periodized
rescaling. Every translate by a vector of length at most \(\rho/2\)
of a point of \(N_\rho(\Gamma)\) belongs to \(N_{2\rho}(\Gamma)\).
This gives the asserted value on \(N_\rho(\Gamma)\), and the mean is
\(|N_{2\rho}(\Gamma)|\).

Parseval gives an \(\ell^2\)-norm \(O(\rho^{1/2})\) for the Fourier
coefficients of the indicator. The bump coefficients satisfy
\(|\widehat\psi(\rho\ell/2)|\ll_M(1+\rho|\ell|)^{-M-2}\).
The sum of \((1+\rho|\ell|)^{-4}\) over the lattice is \(O(\rho^{-2})\):
the ball of radius \(\rho^{-1}\) and the shells with radii
\(2^j/\rho\) each contain \(O((2^j/\rho)^2)\) points. Thus the
\(\ell^2\)-norm of the bump coefficients over \(|\ell|>X\) is
\(O_M(\rho^{-1}(1+\rho X)^{-M})\). Cauchy--Schwarz applied to the
product of the two sets of coefficients proves the tail bound.
Including the mean gives the full sum.
\end{proof}

\begin{proposition}[Zero mass on the swept limsup]\label{prop:curve-upper}
Let \(\tau>1\) and let \(\gamma\) have finite length. For a finite
positive measure \(\mu\) on the torus with \(\widehat\mu(k)=O(H(k))\),
there is \(\nu>0\) such that, for all sufficiently large \(n\),
\begin{equation}\label{eq:curve-upper}
               \mu(U_n(I))\ll h(n)^{2/3}+e^{-\nu n}.
\end{equation}
Consequently \(\mu(V)=0\) if \(\sum_{n\in S}h(n)^{2/3}<\infty\).
\end{proposition}
\begin{proof}
Put \(\Gamma=\gamma(I)\). Since \(\|A^{-n}\|\ll\lambda^n\), the
pullback \(A^{-n}U_n(I)\) lies in a neighborhood of \(\Gamma\) of
width \(O(c\lambda^{-(\tau-1)n})\). With \(\rho=h(n)^{2/3}\), this
width is less than \(\rho\) for large \(n\), by the polynomial lower
bound on \(h\). Test against \(F_\rho\circ A^{-n}\).
The automorphism makes \((A^{-n})^{\mathsf T}\ell\) integral. For
\(0<|\ell|\le\lambda^{n/2}\), its stable dual coordinate and
\eqref{eq:norm} give
\[
  |(A^{-n})^{\mathsf T}\ell|
       \gg \lambda^n|\ell_s|
       \gg \lambda^n/|\ell|
       \ge \lambda^{n/2}.
\]
Thus \(L((A^{-n})^{\mathsf T}\ell)\gg n\), and monotonicity and doubling
give a bound \(O(h(n))\) for the corresponding coefficients of \(\mu\).
The Fourier series of the test function is absolutely convergent by
\eqref{eq:curve-majorant}. The constant term contributes \(O(\rho)\),
and the nonzero terms up to the cutoff contribute
\(O(h(n)\rho^{-1/2})\). Above the cutoff use the mass of \(\mu\) to
bound its coefficients. The tail in \eqref{eq:curve-majorant} there is
exponentially small, since \(\rho^{-1}\) grows at most polynomially in
\(n\). Substituting \(\rho=h(n)^{2/3}\) yields \eqref{eq:curve-upper}.
Borel--Cantelli gives the claim on \(V\).
\end{proof}

\subsection{Graph estimates and the unthinned construction}
Assume now that \(\gamma\) is regular and \(C^2\) and has a point of
nonzero curvature in every open parameter interval. Fix a basis
interval \(J\). In a lift, curvature is nonzero on some open subinterval
of \(J\). On that subinterval, a stable tangency
\(\widetilde\gamma'_u=0\) has \(\widetilde\gamma''_u\ne0\), by the
nonvanishing determinant of the first two derivatives. Such a tangency
is isolated there. We can therefore choose a smaller graph portion
\(u\mapsto u e_u+f(u)e_s\), on which \(|f''|\) is bounded below. This
follows also from the determinant of the two derivatives of the graph.
An unstable tangency corresponds to \(f'=0\) and is allowed on this
portion. Take a nonnegative smooth density \(g(u)\) of integral one,
compactly supported in the graph interval, and let \(\sigma\) be its
image measure. For the graph transform we use the lift, so at real
\(\ell\),
\[
 \widehat\sigma(\ell)=\int g(u)e(-\ell_u u-\ell_s f(u))\,du .
\]
At integer frequencies this agrees with the transform on the torus.

\begin{lemma}[Oscillation on a curved graph]\label{lem:curve-oscillation}
There is a constant \(C\) such that
\begin{equation}\label{eq:curve-oscillation}
|\widehat\sigma(\ell)|\ll
\begin{cases}
 (1+|\ell_s|)^{-1/2},&|\ell_u|\le C(1+|\ell_s|),\\
 (1+|\ell_u|)^{-1},&|\ell_u|> C(1+|\ell_s|).
\end{cases}
\end{equation}
In particular, \(|\widehat\sigma(\ell)|\ll(1+|\ell|)^{-1/2}\).
\end{lemma}
\begin{proof}
Let \(\phi(u)=\ell_u u+\ell_s f(u)\). For \(|\ell_s|\ge1\),
\(\phi'\) is monotone and \(|\phi''|\gg|\ell_s|\). In the classical
second-derivative argument \cite{VdC,Stein}, the set where
\(|\phi'|\le|\ell_s|^{1/2}\) has length \(O(|\ell_s|^{-1/2})\).
On each of its at most two complementary intervals, integrate by parts
using \(1/\phi'\). The boundary terms and the term with \(g'\) are
\(O(|\ell_s|^{-1/2})\). The term with the derivative of \(1/\phi'\)
has the same bound, since its integral in absolute value is the
variation of the monotone function \(1/\phi'\) on that interval.
For \(|\ell_s|<1\) the first estimate follows from the mass of \(\sigma\).

In the second region choose \(C>2\|f'\|_\infty+2\). Then
\(|\phi'|\ge|\ell_u|/2\) on the support interval, and integration by
parts on the full interval bounds the transform by
\(O(|\ell_u|^{-1}+|\ell_s||\ell_u|^{-2})=O(|\ell_u|^{-1})\).
Equivalence of the Euclidean and dual-coordinate norms gives the radial
estimate in both regions.
\end{proof}

Let \(G_n=\chi_{r_n}*((A^n)_*\sigma)\) as in \eqref{eq:tube}.
It is a smooth probability density compactly supported in \(U_n(J)\),
since the parameter density is compactly supported in the chosen portion
of \(J\) and the smoothing radius is at most \(r_n/4\). The ball bound
of Lemma~\ref{lem:ball} gives
\begin{equation}\label{eq:curve-balls}
   \int_{B(y,R)}G_n\ll R^{d_0}\qquad(0<R\le1,\ \tau\ge1).
\end{equation}
It is preserved by convolution with any probability, by integrating
the bound over translates of the ball. The original graph measure
satisfies \(\sigma(B(y,R))\ll R\) for \(0<R\le1\). Indeed only boundedly
many lifts of such a ball can meet its bounded lifted graph, and each
ball projects to an interval of length \(O(R)\) in the \(u\)-coordinate.

Here is the energy consequence of \eqref{eq:curve-balls} that we will
use for both positive densities and signed errors. If a positive
density \(F\) of bounded mass satisfies the ball bound and
\(|f_0|\le qF\), then \(\|f_0\|_s\ll_s q\) for \(0<s<d_0\), with the
norms of Section~\ref{sec:construction}. In fact its positive heat
kernel representation bounds \(\|f_0\|_s^2\) by a constant times
\(q^2\iint (1+\operatorname{dist}(x,y)^{-s})F(x)F(y)\,dx\,dy\).
The integral is bounded by the ball estimate, since the singular part
is controlled by \(\int_0^1 R^{d_0-s-1}\,dR\).

For the unthinned lower bound, fix \(k\ne0\) and put
\[
 t_k=\frac{\log |k_s/k_u|}{2\log\lambda},\qquad
 \max(\lambda^n|k_u|,\lambda^{-n}|k_s|)
       =\sqrt{|k_u k_s|}\,\lambda^{|n-t_k|}.
\]
Equations \eqref{eq:norm}, \eqref{eq:tube}, and the radial estimate
give \(\sum_n|\widehat G_n(k)|\ll1\) uniformly in \(k\ne0\), by
summing \(\lambda^{-|n-t_k|/2}\). To obtain the weighted bound
\begin{equation}\label{eq:curve-weighted}
                  \sum_n h(n)|\widehat G_n(k)|\ll H(k),
                  \qquad k\ne0,
\end{equation}
write \(L=L(k)\) and fix \(\kappa>0\) small and \(C_2>0\) large.
If \(n\le\kappa L\), with \(\kappa\tau\log\lambda<1/2\), then
\(r_n|k|\gg_c e^{L/2}\) for large \(L\). The Schwartz bound for the
bump gives \(O(e^{-\nu L})\) per term for some \(\nu>0\) across
these \(O(L)\) times. If \(n\ge C_2 L\), with \(C_2\log\lambda>2\),
then \(\lambda^n|k_u|\gg \lambda^n/|k|\gg e^{n\log\lambda/2}\).
The radial graph estimate gives an exponentially summable tail.
The polynomial lower bound on \(H(k)\) absorbs both contributions.
On the intervening range \(h(n)\asymp H(k)\), so the uniform unweighted
sum suffices. The bounded range of \(L\) follows from the same
unweighted sum and positivity of \(H\) at each of the finitely many
corresponding lattice points.

Suppose \(\sum_{n\in S}h(n)=\infty\). In any time tail, normalize a
finite sum of the densities \(h(n)G_n\) by its total weight to obtain
\(Q\). Its support is compact in the union of the corresponding
\(U_n(J)\), and \eqref{eq:curve-weighted} gives
\(\|Q-1\|_H\ll(\sum h(n))^{-1}\). To see convergence in the second norm
of Lemma~\ref{lem:compact}, fix \(0<s<s'<d_0\). The ball bound for
the average gives a uniform \(s'\)-norm bound. Above frequency \(X\),
the \(s\)-norm of \(Q-1\) is then \(O((1+X)^{-(s'-s)/2})\).
At bounded frequencies the \(H\)-norm controls the \(s\)-norm of
the error. Taking \(X\) large and then the total weight large makes
both errors small, as in Lemma~\ref{lem:average}. Lemma~\ref{lem:compact}
gives the compact probability on \(D\) with decay \(o(H)\) at real
frequencies and all subcritical energies.

\subsection{A square sum over a thin frequency strip}
\begin{lemma}[Thin strips in dual coordinates]\label{lem:curve-strip}
For \(B\in\mathbb R\), \(0\le w\le1\), and \(0<p\le1\), set
\[
 W(B,w,p)^2
  =\sum_{\substack{0\ne\ell\in\mathbb Z^2\\|\ell_s-B|\le w}}
      |\widehat\sigma(\ell)|^2\min\{1,p|\ell|\}^2.
\]
Then
\begin{equation}\label{eq:curve-strip}
 W(B,w,p)\ll \sqrt w+\sqrt p
 \begin{cases}
      |B|^{1/2}, & |B|<1,\\
      \min\{(p|B|)^{1/2},(p|B|)^{-1/2}\}, & |B|\ge1 .
 \end{cases}
\end{equation}
For all sufficiently small \(|B|+w\), one also has
\begin{equation}\label{eq:curve-strip-small}
                  W(B,w,p)\ll |B|+w .
\end{equation}
\end{lemma}
\begin{proof}
For \(w>0\), applying \eqref{eq:norm} to differences of distinct
lattice points in the strip gives separation \(\gg w^{-1}\) of their
unstable coordinates. A zero-width strip contains at most one lattice
point. Put \(X=1+|B|\), so \(1+|\ell_s|\asymp X\) on the strip.
In \(|\ell_u|\le C'X\) there are \(O(1+wX)\) points, and
\eqref{eq:curve-oscillation} bounds each transform by \(O(X^{-1/2})\).
Since \(|\ell|\ll X\) in this region, its contribution to \(W^2\) is
\[
               O\bigl(\min(1,pX)^2(X^{-1}+w)\bigr).
\]
Choose \(C'\) large enough that \(|\ell_u|>C'X\) is in the nonstationary
region of \eqref{eq:curve-oscillation}. There \(|\ell|\asymp|\ell_u|\),
and the compensated transform is \(O(\min\{p,|\ell_u|^{-1}\})\).
On either side of the origin a sum of a decreasing function of
\(|\ell_u|\) is at most its first possible value plus \(O(w)\) times
its tail integral, by the separation. For the squares this yields
\[
 O\bigl(\min\{p^2,X^{-2}\}+w\min\{p,X^{-1}\}\bigr),
\]
because \(\int_X^\infty\min\{p^2,u^{-2}\}\,du\le2\min\{p,X^{-1}\}\).
The argument with just the first possible value also covers \(w=0\).

For \(x=|B|+w\) small, \eqref{eq:norm} instead gives
\(|\ell_u|\gg x^{-1}\) at every nonzero point. These points are all
nonstationary. Summing \(O(|\ell_u|^{-2})\) without the compensating
factor gives \(W^2\ll x^2+wx\ll x^2\), with an empty sum when \(x=0\).
This proves \eqref{eq:curve-strip-small}. For \(|B|\ge1\), the square
roots of the preceding general bounds are
\(O(\sqrt w+X^{-1/2}\min\{1,pX\})\), giving the second profile in
\eqref{eq:curve-strip}. For \(|B|<1\), the general bound is
\(O(p+\sqrt w)\). Together with \eqref{eq:curve-strip-small} it gives
\(O(\sqrt w+\min\{p,x\})\), also when \(x\) is bounded below since
\(p\le1\). Finally \(\min\{p,x\}\le\sqrt{px}\le\sqrt p\,|B|^{1/2}
+\sqrt w\), giving the first profile.
\end{proof}

\subsection{Compensated multipliers on separated schedules}
Take a nonnegative Schwartz probability \(\eta\) on \(\mathbb R^2\)
with compactly supported Fourier transform. For example, normalize the
squared modulus of the inverse transform of a nonzero smooth
compactly supported function. The transform of this square is a
convolution of two compactly supported functions. For \(0<p\le1\), let
\(\eta_p\) be the periodized rescaling and set
\[
 D_{n,p}=G_n*((A^n)_*\eta_p),\qquad
 \widehat D_{n,p}(k)=\widehat G_n(k)\,
                         \widehat\eta(p(A^n)^{\mathsf T}k).
\]
These are smooth probability densities with the ball bound
\eqref{eq:curve-balls}. The density of the pullback convolution
\(\sigma*\eta_p\) is \(O(p^{-1})\). Indeed, for \(M>2\), the
periodized kernel is
\(O_M(p^{-2}(1+\operatorname{dist}(z,0)/p)^{-M})\).
To see this, keep the nearest lattice translates and sum the remaining
Schwartz bounds \(O_M(p^{M-2}(1+|m|)^{-M})\) over lattice translates
\(m\). Annuli of radii \(2^j p\), together with
\(\sigma(B(y,R))\ll R\), bound the convolution by
\(O(p^{-1}\sum_{j\ge0}2^{j(1-M)})\). Pushing forward by \(A^n\)
preserves the density bound, and convolution with \(\chi_{r_n}\)
preserves it again. Thus \(\|D_{n,p}\|_\infty\le C_0/p\).
The Fourier support is contained in \(|(A^n)^{\mathsf T}k|\ll p^{-1}\).
In particular,
\begin{equation}\label{eq:curve-band}
 |k|\ll\lambda^n/p,\qquad
 |k_s|\gg\lambda^n p \quad\text{if }k\ne0
 \text{ and }\widehat D_{n,p}(k)\ne0.
\end{equation}
The second bound follows from \(|k_u|\ll\lambda^{-n}/p\) and
\eqref{eq:norm}. The finite first moment of \(\eta\) gives
\begin{equation}\label{eq:curve-cancel}
                 |1-\widehat\eta(p\ell)|\ll\min\{1,p|\ell|\}.
\end{equation}

At a single level, \(1+c_0p(G_n-D_{n,p})\) is a probability density
for \(c_0 C_0\le 1/2\), and its mass on \(U_n(J)\) is at least \(c_0p\).
Its nonzero coefficients are \(O(p^{3/2})\), since
\(\min\{1,p|\ell|\}(1+|\ell|)^{-1/2}\le\sqrt p\) in the radial
estimate. The following construction allocates these corrections to
successive levels of a finite multiplier.

\begin{proposition}[Positive multipliers from logarithmic gaps]
\label{prop:curve-compensation}
Let \(\tau\ge1\). Suppose that for every \(K>0\) there is
\(T\subseteq S\) with \(\sum_{n\in T}h(n)^{2/3}=\infty\) and with
successive gaps \(m-n\ge K\log(e+n)\).
Then the hypotheses of Lemma~\ref{lem:compact} hold for \(d_0=1+1/\tau\).
\end{proposition}
\begin{proof}
Fix \(J\), \(0<s<s'<d_0\), and a positive scalar \(v\). Choose
\(p_0\in(0,1)\) and \(c_0>0\) so that \(c_0 C_0\le1/2\) and
\(\theta=(1-c_0p_0)^{-1/3}<\lambda\).
We will take a large gap constant \(K\) and use \(T\) starting at a
late threshold \(N\). At a trial time \(n\), let \(R_{<n}\) be the
product of the factors \(1-c_0p_m D_{m,p_m}\) at earlier used times
\(m<n\), starting with the constant one. Set
\begin{equation}\label{eq:curve-adapt}
        M_n=\int R_{<n},\qquad
        p_n=v h(n)^{2/3}M_n^{-1/3}.
\end{equation}
Stop before using the first trial with \(p_n>p_0\). At all used
levels write \(D_n=D_{n,p_n}\). Each factor lies between \(1/2\)
and one, so \(M_n>0\), \(0\le R_{<n}\le1\), and
\(\|R_{<n}\|_2\le\sqrt{M_n}\) at every trial. In particular
\(p_n\ge v h(n)^{2/3}\).

\medskip
\noindent\emph{Separation and exact mass.}
By \eqref{eq:curve-band}, both the norm and the absolute stable
coordinate of every index in the Fourier support of \(R_{<n}\) are
at most
\[
        Y_n=C_3\sum_{\substack{m<n\\m\ {\rm used}}}\frac{\lambda^m}{p_m},
        \qquad w_n=\lambda^{-n}Y_n ,
\]
with a fixed sufficiently large \(C_3\). For any prescribed
\(L_0>2\beta/3\), we can arrange
\begin{equation}\label{eq:curve-width}
                         w_n\le n^{-L_0}
\end{equation}
at every trial, including the stopping trial. In fact, if \(m_*\) is
the preceding time, monotonicity and the geometric sum give
\(Y_n\ll v^{-1}h(n)^{-2/3}\lambda^{m_*}\).
If \(m_*\ge n/2\), the gap gives
\(\lambda^{m_*-n}\le(e+n/2)^{-K\log\lambda}\).
If \(m_*<n/2\), the same factor is exponentially small in \(n\).
With no preceding time the support bound is zero. The polynomial
bound \(h(n)^{-2/3}\ll n^{2\beta/3}\) therefore gives
\eqref{eq:curve-width} by choosing \(K\) large and then \(N\) large
for fixed \(v\). We may increase \(L_0\) before choosing \(K\) and
\(N\), as needed in the estimates below.

At a used level, \(w_n\) is smaller than the fixed constant in the
lower bound \(|k_s|\gg\lambda^n p_n\) of \eqref{eq:curve-band}, times
\(p_n\), by \eqref{eq:curve-adapt} and \eqref{eq:curve-width}.
Thus the Fourier pairing \(\int R_{<n}D_n\) has just the constant
term \(M_n\). For the next trial \(n'\) this gives
\[
          \int R_{<n}D_n=M_n,\qquad M_{n'}=(1-c_0p_n)M_n .
\]
Concavity of the cube root shows
\(M_{n'}^{1/3}\le M_n^{1/3}-(c_0v/3)h(n)^{2/3}\).
If there were no stop, divergence on \(T\) would contradict positivity
of these masses. Write \(n_*\) for the stopping time. Its defining
inequality gives
\begin{equation}\label{eq:curve-residual-mass}
                   M_{n_*}<(v/p_0)^3 h(n_*)^2.
\end{equation}
The starting threshold can also be chosen so that the first trial
satisfies \(p_n\le p_0\).

Set
\[
   Q=\sum_{n\ {\rm used}} c_0p_n R_{<n}G_n,\qquad
   E_n=c_0p_n R_{<n}(G_n-D_n).
\]
This \(Q\) is smooth, nonnegative, and supported compactly in the
finite union of the selected \(U_n(J)\). Telescoping the successive
differences of \(R_{<n}\) gives
\begin{equation}\label{eq:curve-telescope}
                       Q-1=\sum_{n\ {\rm used}} E_n-R_{<n_*}.
\end{equation}
On the residual Fourier support, \(|k|\le Y_{n_*}\le\lambda^{n_*}\),
so \(H(k)\gg h(n_*)\) by gauge regularity. Nonnegativity of the
residual bounds each coefficient by \(M_{n_*}\). Its \(s\)-norm
is at most its \(L^2\)-norm since \(s<2\). In particular
\begin{equation}\label{eq:curve-residual}
 \|R_{<n_*}\|_H\ll v^3 h(n_*),\qquad
 \|R_{<n_*}\|_s\le\sqrt{M_{n_*}}\ll v^{3/2}h(n_*).
\end{equation}

\medskip
\noindent\emph{Weighted Fourier bound.}
In the Fourier convolution for \(E_n\), let \(l\) be an index of
\(R_{<n}\) and put \(\ell=(A^n)^{\mathsf T}(k-l)\) and
\(B_n=b^n k_s\). These \(\ell\) are distinct integer points with
\(|\ell_s-B_n|\le w_n\). The term \(\ell=0\) vanishes by compensation.
Using \eqref{eq:curve-cancel}, Cauchy--Schwarz, and Parseval gives
\begin{equation}\label{eq:curve-error-coeff}
 |\widehat E_n(k)|\ll p_n\sqrt{M_n}\,W(B_n,w_n,p_n)
       \sup_{|l|\le Y_n}|\widehat\chi(r_n(k-l))|.
\end{equation}
Here \(\sqrt{M_n}\) bounds the \(\ell^2\)-norm of the coefficients of
\(R_{<n}\). The bump supremum is always at most one.

Use \(L=L(k)\) and the constants \(\kappa,C_2\) from the weighted
graph sum, enlarging \(C_2\) if needed. For \(n\le\kappa L\),
\(Y_n\le\lambda^n\) implies \(|k-l|\ge|k|/2\) for all relevant \(l\)
and all large \(L\). Then \(r_n|k-l|\gg_c e^{L/2}\), and the bump
supremum is \(O(e^{-\nu L})\) for some \(\nu>0\). The strip bound
is uniform, so the sum in this range divided by \(H(k)\) is
\(O(L^{1+\beta}e^{-\nu L})\). A nonempty early range has
\(L\ge N/\kappa\), hence this tends to zero uniformly as \(N\) increases.

For \(n\ge C_2 L\), \(|B_n|\le |k|\lambda^{-n}\ll e^{-\nu' n}\)
for some \(\nu'>0\). The small-strip estimate
\eqref{eq:curve-strip-small} gives
\(W(B_n,w_n,p_n)\ll e^{-\nu'n}+n^{-L_0}\). Since \(H(k)^{-1}\ll n^\beta\)
in this range, the sum divided by \(H(k)\) is bounded by the tail of
\(\sum_n n^\beta(e^{-\nu'n}+n^{-L_0})\), convergent for
\(L_0>\beta+1\).

In the middle range \(\kappa L<n<C_2 L\), \(h(n)\asymp H(k)\).
The \(\sqrt{w_n}\) term in \eqref{eq:curve-strip}, after division by
\(H(k)\), contributes at most a constant times the tail of
\(\sum_n n^{\beta-L_0/2}\). For the remaining terms,
\[
                   p_n^{3/2}\sqrt{M_n}=v^{3/2}h(n)
\]
by \eqref{eq:curve-adapt}. The profiles accompanying this factor have
bounded sum over the used times, uniformly in \(k\). For \(|B_n|<1\),
the sum of \(|B_n|^{1/2}\) is geometric because \(|B_n|\) decreases
by at least the factor \(\lambda^{-1}\) at each successive time.
For the other profile let \(z_n=p_n|B_n|\). At consecutive used
times \(n<n'\), monotonicity of \(h\) and the exact mass recurrence
give \(p_{n'}\le\theta p_n\), hence \(z_{n'}\le(\theta/\lambda)z_n\).
The sum of \(\min\{z_n^{1/2},z_n^{-1/2}\}\) on \(|B_n|\ge1\) is
bounded by the geometric sums before and after \(z_n\) crosses one.
Taking \(L_0>2(\beta+1)\) in these three ranges gives
\begin{equation}\label{eq:curve-error-H}
              \left\|\sum_{n\ {\rm used}}E_n\right\|_H
                 \ll v^{3/2}+o_{N\to\infty}(1).
\end{equation}
The convergence of the error is uniform in the stopping time for
fixed \(v\).

\medskip
\noindent\emph{Energy bound.}
The pointwise inequality \(|E_n|\le c_0p_n(G_n+D_n)\) and the
positive kernel estimate following \eqref{eq:curve-balls} give
\(\|E_n\|_{s'}\ll p_n\). Fix \(J_0>0\). On \(|k|>n^{J_0}\), the
comparison of the two Fourier weights bounds the \(s\)-norm of these
modes by \(O(p_n n^{-J_0(s'-s)/2})\). For \(|k|\le n^{J_0}\),
\eqref{eq:curve-strip-small} and \eqref{eq:curve-error-coeff} give
\[
           |\widehat E_n(k)|\ll
              p_n(\lambda^{-n}n^{J_0}+n^{-L_0}).
\]
The sum of \((1+|k|)^{s-2}\) over these \(k\) is \(O(n^{J_0s})\)
by dyadic lattice counting. The \(s\)-norm of the low modes is
therefore \(O(p_n n^{J_0s/2}(\lambda^{-n}n^{J_0}+n^{-L_0}))\).
Choose \(J_0(s'-s)/2>2\) and then \(L_0>J_0s/2+2\), also meeting
the earlier requirements. With \(p_n\le p_0\), both norm bounds
are summable over \(n\), and their tails tend to zero uniformly in
the stopping time.

First choose \(v\) small, then take \(N\) large in the subsequence
corresponding to the required \(K\). Equations
\eqref{eq:curve-telescope}--\eqref{eq:curve-error-H} and the energy
bound make \(\|Q-1\|_H+\|Q-1\|_s\) as small as desired. The choice of
\(N\) can place the finite support union inside any specified time
tail. This proves the proposition.
\end{proof}

\subsection{The logarithmic endpoint and a mixed arc}
\begin{corollary}[Thinning a schedule]\label{cor:curve-log-gaps}
For \(\tau\ge1\), the compact-support conclusion on \(D\) with decay
\(o(H)\) and all \(0<s<d_0\) energies holds whenever
\[
                   \sum_{n\in S}\frac{h(n)^{2/3}}{\log(e+n)}=\infty .
\]
\end{corollary}
\begin{proof}
For fixed \(K>0\), choose successive times greedily, taking after \(n\)
the first available time at least \(n+\lceil K\log(e+n)\rceil\).
Every group consisting of a chosen \(n\) and the skipped times before
the next choice has \(O_K(\log(e+n))\) elements. For every \(m\) in the
group one has
\(h(m)^{2/3}/\log(e+m)\le h(n)^{2/3}/\log(e+n)\). Thus the group
contributes \(O_K(h(n)^{2/3})\) to the divergent sum. The chosen
times retain divergence of \(\sum h(n)^{2/3}\).
Proposition~\ref{prop:curve-compensation} and Lemma~\ref{lem:compact}
give the conclusion. The latter uses nested compact supports in fresh
time tails, so each point of the limit support belongs to the
intersection in \eqref{eq:dense}. It supplies all the subcritical
energies on the same measure and transfers the decay to real
frequencies in its local lift.
\end{proof}

For \(S=\mathbb N\), take \(h(t)=t^{-3/2}\). The thinning sum is
the divergent harmonic logarithmic series. For \(\tau>1\) and
\(q>3/2\), the gauge \(h(t)=t^{-q}\) gives a convergent obstruction
sum in Proposition~\ref{prop:curve-upper}. This proves the sharp
power in Theorem~\ref{thm:planar}, including its little-\(o\) endpoint.
For schedules whose successive gaps divided by \(\log(e+n)\) tend
to infinity, every prescribed gap condition holds on a tail of \(S\).
Consequently convergence or divergence of \(\sum_{n\in S}h(n)^{2/3}\)
decides the corresponding alternatives of Theorem~\ref{thm:sparse}
on these schedules. The estimates leave open the full-time
refinements \(h(t)=t^{-3/2}(\log(e+t))^{-q}\) for \(0<q\le3/2\),
as well as the general case where thinning loses divergence.

For any infinite \(S\), choose breakpoints \(N_j\) with
\(N_{j+1}>2N_j\) and \(S\cap[N_j,N_{j+1})\ne\varnothing\).
The gauge equal to \(1/j\) on each of these intervals, and to one
before the first, is regular since doubling crosses at most one
breakpoint. Its sum over \(S\) diverges. The unthinned construction
gives local Hausdorff dimension at least \(d_0\) for \(D\) by finite
energy. The length cover in Section~\ref{sec:construction} bounds the
dimension of \(V\) by \(d_0\), and the pullback cover there gives
\(\dim_H E_y=0\) for all \(y\) when \(\tau>1\), with the stated
convention for empty fibers. Proposition~\ref{prop:curve-upper}
excludes any positive power Fourier decay for a measure carried by
\(V\), since it implies the bound for every \(h(t)=t^{-q}\).
These observations give the asserted dimensions for \(D,P,V\).

\begin{example}[A mixed straight portion joined to a curved portion]
\label{ex:curve-mixed}
Let \(S=\mathbb N\), \(\tau>1\), and \(\varepsilon>0\). Use the lift
\[
 t\longmapsto \eta\left[t e_u+
        \bigl(t+\varepsilon(t-\tfrac12)_+^3\bigr)e_s\right],
 \qquad 0\le t\le1,\qquad x_+=\max\{x,0\}.
\]
Choose \(\eta>0\) so that the lifted graph lies in a ball of radius
less than \(1/4\). It projects to an embedded \(C^2\) arc with mixed
straight first half and nonzero curvature on \((1/2,1)\).

Fix a closed interval \(J\) of positive length inside the straight
half. Every dense fiber meets \(J\), so \(D\) lies in the swept
limsup for the mixed segment \(\gamma(J)\). By Theorem~\ref{thm:sparse},
every finite positive torus measure with decay \(O(L^{-q})\), \(q>1\), gives that limsup
zero mass. Conversely, every open parameter interval contains a smaller
curved or mixed straight portion. The curved portion gives
\eqref{eq:curve-weighted}; on the straight portion the mixed-line
estimate in Section~\ref{sec:lines} gives the same bound with
\(w_n=h(n)\). Both have the ball bound of Lemma~\ref{lem:ball}.
For \(h(t)=t^{-1}\) the weights sum to infinity. Cycling these portions
in Lemmas~\ref{lem:average} and~\ref{lem:compact} gives \(o(L^{-1})\)
on compact dense-fiber support.

Applying the curved construction on a compact subinterval of
\((1/2,1)\) gives power \(3/2\) on a compact subset of \(P\), since
its dense-fiber centers have a nonempty fiber for the full arc.
Proposition~\ref{prop:curve-upper} for the entire arc excludes all
larger powers already on \(V\). Both constructions give all the
subcritical energies, so the distinction between the sharp powers
persists in every target patch at the same local Hausdorff dimension.
\end{example}

\subsection{The radius \texorpdfstring{\(\tau=1\)}{tau = 1}}
For the curved arcs of Theorem~\ref{thm:planar}, at \(\tau=1\) the
sets \(D,P,V\) are Haar-conull along any fixed infinite \(S\). To see
this, in each basis interval \(J\) take a transverse graph portion.
Its translates by stable displacements \(u e_s\) with \(|u|<c/2\)
form a nonempty open strip in the pullback torus, using an open
graph interval. Call its projection \(O_J\). If \(A^{-n}y\in O_J\),
then \(y\in U_n(J)\), since the displacement after applying \(A^n\)
has Euclidean length less than \((c/2)\lambda^{-n}\).

Write \(m\) for Haar probability. The nonzero dual orbits of \(A^{-1}\)
grow in norm by \eqref{eq:norm}, so \(A^{-1}\) is mixing. Indeed, for
trigonometric polynomials the only correlations that persist are
those of the constant terms, and \(L^2\)-approximation extends this
to all \(L^2(m)\) functions. Enumerate \(S=\{n_1<n_2<\cdots\}\).
The averages of \(\mathbf 1_{O_J}(A^{-n_j}y)\) over the first \(N\)
times converge in \(L^2(m)\) to \(m(O_J)>0\): for every tolerance
\(\epsilon>0\), mixing bounds the covariances in the variance sum
with \(|n_i-n_j|\ge K\) by \(\epsilon\) for large enough \(K\).
The remaining ordered pairs number \(O(KN)\) because the \(n_i\)
are distinct integers. Dividing by \(N^2\) and taking \(N\) large
gives convergence. An almost everywhere convergent subsequence of
the averages then implies infinitely many visits for almost every
center. Intersecting over the countable basis yields the claim by
\eqref{eq:dense}.

Every smooth probability density in a target patch therefore gives
\(D\) full mass. Inner regularity gives a compact subset of \(D\) of
positive Haar measure in each patch. Its normalized Haar restriction
has bounded density and every \(0<s<2\) energy. The multiplier proofs
also give their compact Fourier lower conclusions at \(\tau=1\),
with \(d_0=2\) by Lemma~\ref{lem:ball}. The optimal decay for compact
supports at this radius remains open here.

\begin{remark}[Fibers at the critical radius]
When the stable tangencies are countable, every fiber at \(\tau=1\)
has Hausdorff dimension zero. The inverse image of a target ball has
stable width \(O(1)\) and unstable width \(O(\lambda^{-2n})\).
It has only \(O(1)\) lattice lifts meeting the bounded lifted arc.
On each parameter chart where \(|\widetilde\gamma'_u|\) is bounded
below, the hits at time \(n\) are covered by \(O(1)\) intervals of
length \(O(\lambda^{-2n})\), by projection onto the unstable axis.
A countable cover by such charts and the countable set of
tangencies therefore give dimension zero.
\end{remark}

\section{A fixed center at the stable threshold}\label{sec:critical}

Throughout this section \(A\in GL_2(\mathbb Z)\) is hyperbolic,
\(\lambda>1\) is its expanding modulus, and \(\chi=\log\lambda\).
Write \(a,b\) for the signed eigenvalues, so that
\(|a|=\lambda\) and \(|b|=\lambda^{-1}\), and \(e_u,e_s\) for unit right
eigenvectors. For \(z=z_u e_u+z_s e_s\), the subscripts denote
coordinates in this basis. Let the lift of the segment be \(p+xv\),
\(x\in I\), where \(I\) is a nondegenerate compact interval.
Transversality means \(v_u\ne0\). Use open eigenrectangles of
half-widths \(c_u e^{-\alpha n}\), \(c_s e^{-\beta n}\), with all
four constants positive. For visits to centers \(y_n\) along \(S\),
write \(H_n\) for the set of hitting parameters at time \(n\) and
\(E=\limsup_{n\in S} H_n\). For fixed centers and balls this is \(E_y\).
Balls of radius \(c e^{-\alpha n}\) contain and are contained in
such rectangles with \(\beta=\alpha\), by equivalence of the
Euclidean and eigen-coordinate norms.

\subsection{Slice dimension and access}

\begin{proposition}\label{prop:transverse-slices}
For every transverse affine segment, every center sequence and every
infinite \(S\), one has
\[
\dim_H E=
\begin{cases}
(\chi-\beta)/(\chi+\alpha),&0<\beta<\chi,\\
0,&\beta\ge\chi .
\end{cases}
\]
We use \(\dim_H\varnothing=0\). In the positive-dimensional case
\(\mathcal H^{(\chi-\beta)/(\chi+\alpha)}(E\cap J)=\infty\) for every
nonempty relatively open subinterval \(J\) of \(I\).
For balls \(B(y_n,c\lambda^{-\tau n})\) the corresponding dimension
is \((1-\tau)/(1+\tau)\) for \(0<\tau<1\), and zero for \(\tau\ge1\).
\end{proposition}

\begin{proof}
Unimodularity gives \(A^{-n}\mathbb Z^2=\mathbb Z^2\). Thus \(H_n\)
is the intersection, in parameter coordinates, with rectangles
about the lifts of \(q_n=A^{-n}y_n\) of half-widths
\begin{equation}\label{eq:inverse-slice-widths}
 u_n=c_u e^{-(\chi+\alpha)n},\qquad W_n=c_s e^{(\chi-\beta)n}.
\end{equation}
For bounded \(W_n\) there are only
boundedly many relevant lifts, uniformly in \(q_n\), since the lifts
lie in a fixed bounded neighborhood of the segment. Each gives a
parameter interval of length at most \(2u_n/|v_u|\).

For \(W_n\to\infty\), any two relevant lifts have stable
separation \(O(W_n)\). For a nonzero lattice vector \(z\), the
integer \(\det[z,Az]=(b-a)\det[e_u,e_s]z_u z_s\) is nonzero,
since a contracting eigenvalue cannot be a nonzero rational
algebraic integer and the characteristic polynomial is irreducible.
Thus \(|z_u z_s|\gg1\), so the unstable coordinates of relevant
lifts are separated by \(\gg W_n^{-1}\). Their unstable range is
bounded, yielding \(O(W_n)\) parameter intervals of diameter
\(O(u_n)\). These covers give both upper bounds.

For the lower bound take a covering radius \(C_0\) of the lattice in
the eigen-coordinate maximum norm. Choose \(m\) so that
\(C_0\lambda^m\le W_n/4\) and \(\lambda^m\asymp W_n\).
Apply \(A^m\) to \(p+xv\) and to the translated lattice of lifts of
\(q_n\). The resulting lattice translate still has covering radius
\(C_0\). Undoing \(A^m\) at a nearest lift gives a center with
unstable discrepancy \(O(W_n^{-1})\) and stable discrepancy at most
\(W_n/4\). Moving the parameter to the exact unstable match changes
the stable coordinate by \(O_v(W_n^{-1})\). On compact intervals in
the interior of \(I\) this yields an \(O(W_n^{-1})\)-net of
parameters around which intervals of radius \(u_n/(2|v_u|)\) lie
in \(H_n\) for large \(n\). The net centers have separation
\(\gg W_n^{-1}\) by the preceding lattice bound.

For \(d=(\chi-\beta)/(\chi+\alpha)\), raising the interval radii to
the power \(d\) makes them comparable to \(W_n^{-1}\). The enlarged
intervals cover a uniform positive fraction of each fixed subinterval
at every sufficiently late \(n\in S\). Apply the net bound inside
the middle of the subinterval and use separation to bound overlap.
The same fraction bounds the measure of their limsup there by
continuity from above for tail unions. The Lebesgue density theorem
then gives a full-measure enlarged limsup. The one-dimensional mass
transference principle gives the claimed infinite Hausdorff measure
for the original intervals \cite[Theorem~2]{BV}. The ball statements
follow by the rectangle inclusions.
\end{proof}

For full times the subcritical ball dimension also follows by stable
saturation from the ambient theorem of Hu, Persson, Wu and Zhang
\cite[Theorem~2.1]{HPWZ}. Let \(F_\eta\) be the ambient limsup for
balls of radius \(e^{-\eta n}\) with the given centers. Fix
\(0<\alpha<\chi\) and a small stable flow box over a transverse
subinterval. Stable shifts in this box change the forward images
by \(O(e^{-\chi n})\). For sufficiently small \(\varepsilon>0\), the
saturation of the slice at radius \(c e^{-\alpha n}\) is contained
in \(F_{\alpha-\varepsilon}\). Conversely, \(F_{\alpha+\varepsilon}\)
in the box is contained in that saturation. In the bi-Lipschitz
product coordinates the saturation has dimension equal to the slice
dimension plus one. The ambient dimension and local lower bound
from the large-intersection conclusion squeeze this sum between
\(2\chi/(\chi+\alpha+\varepsilon)\) and
\(2\chi/(\chi+\alpha-\varepsilon)\). Letting \(\varepsilon\downarrow0\)
gives the formula. The cover and net above supply the sparse-time
version.

At \(\beta=\chi\), for a fixed center \(y\), put \(q_n=A^{-n}y\).
If \(q_n\) is dense along \(S\), the critical fiber on every
transverse segment is dense \(G_\delta\). Indeed, each open parameter
interval \(J\) gives a nonempty open set
\begin{equation}\label{eq:critical-open-box}
 \{p+xv+t e_s:x\in J,\ |t|<c_s/2\}\pmod{\mathbb Z^2}.
\end{equation}
A visit of \(q_n\) to this set gives a hit with exact unstable
match and strict stable margin. Removing finitely many points from
a dense orbit sample in the torus leaves it dense. Every parameter
tail union is consequently open and dense, and Baire's theorem
applies. An inscribed critical rectangle gives the same assertion
for balls.

For any one prescribed \(S=\{n_1<n_2<\cdots\}\), Haar-almost every
\(y\) has this density property. For \(T=A^{-1}\), every nonzero
dual integer frequency escapes under
\((T^{\mathsf T})^n\), which proves mixing first for trigonometric
polynomials and then in \(L^2\). For an open basis set \(O\), the
variance of the mean of \(\mathbf 1_O(T^{n_j}y)\), \(j\le m\),
tends to zero. For any fixed \(K\), there are \(O(Km)\) pairs with
gap at most \(K\), and the covariances for larger gaps tend uniformly
to zero as \(K\to\infty\). A subsequence of the means thus converges
almost surely to the positive Haar measure of \(O\), and a countable
basis gives density. The same argument applies to forward samples
of \(A\).

In the dense case the ambient critical limsup has dimension one.
Its inverse rectangles are covered at each time by \(O(u_n^{-1})\)
balls of radius \(O(u_n)\). A hit infinitely often for a narrower
stable half-width gives a short stable segment hitting infinitely
often for the original width, since a stable displacement \(t e_s\)
becomes \(b^n t e_s\). The same proof works for balls by sandwiching.
This is consistent with the critical dependence on the centers in
\cite[Remark~2.2]{HPWZ}.

If \(y\) is periodic, retain the residues of
its period that occur infinitely often in \(S\). At \(\beta=\chi\),
the fiber consists exactly of the finitely many parameters with an
exact unstable match to a lift of a retained orbit point and stable
discrepancy strictly less than \(c_s\). For critical balls the exact
stable threshold is \(c\), since \(e_s\) is a unit vector. For
\(\beta>\chi\) both coordinates must match. To see necessity, pass
to one residue and one bounded lift along infinitely many hits, and
use \eqref{eq:inverse-slice-widths}. After an exact unstable match,
equality at the critical stable boundary fails the open test.
Sufficiency follows at every sufficiently late time of the retained
residue.

For general \(y\), a critical hitting parameter must satisfy
\(p+xv\in\omega_S(q)+[-c_s,c_s]e_s\pmod{\mathbb Z^2}\), where
\(\omega_S(q)\) is the set of subsequential limits along \(S\).
The rate in the unstable coordinate is essential. For example, let
\(y=y_*+t e_u\pmod{\mathbb Z^2}\) with \(y_*\) periodic and \(t\ne0\).
Infinitely many hits would force an exact unstable match to a fixed
lift of \(A^{-n}y_*\) on a subsequence of one residue. On that
subsequence the actual discrepancy is \(|t|\lambda^{-n}\),
whereas \eqref{eq:inverse-slice-widths} requires the strictly smaller
scale \(u_n\). These critical fibers are empty.

\subsection{One target for a prescribed schedule and envelope}

The following bound applies to arbitrary center sequences, with
\(\|\cdot\|_{\mathbb R/\mathbb Z}\) denoting distance to the nearest
integer.

\begin{lemma}\label{lem:parameter-block}
Let \(N\le K\le(1+\delta)N\) and
\(Q_K=(A^Kv)_i\), where \(i\) is an ordinary coordinate for which
\((e_u)_i\ne0\). There is a \(\delta>0\) depending on the exponents
such that, for some \(\eta>0\) and all sufficiently large \(N\),
\[
 \mu\left(\bigcup_{n\in S\cap[N,K]}H_n\right)
 \ll e^{-\eta N}\mu(\mathbb R)
      +\sup_{\ell\in\mathbb Z\setminus\{0\}}
             |\widehat\mu(\ell Q_K)|
\]
for every finite positive parameter measure \(\mu\), uniformly in
the centers.
\end{lemma}

\begin{proof}
Set \(\gamma=\min(\alpha,\beta)\). A hit is a torus error of size
\(O(e^{-\gamma n})\). Propagate the whole error by the integer map
\(A^{K-n}\) and take coordinate \(i\). The result is
\[
 \|Q_Kx-(A^{K-n}y_n-A^Kp)_i\|_{\mathbb R/\mathbb Z}
 \ll e^{-(\gamma-\chi\delta)N}.
\]
There are at most \(\lfloor\delta N\rfloor+1\) phases. Choose
\(\delta\) small and \(\eta_0\) with
\(\delta\log16<\eta_0<\gamma-\chi\delta\).
Apply Lemma~\ref{lem:majorant} to \(Q_Kx\) with
\(M=\lfloor e^{\eta_0 N}\rfloor\). Its constant term gives the
exponential error, and its other coefficients bound the integral by
the displayed Fourier supremum. Transversality ensures
\(|Q_K|\asymp_v\lambda^K\) for large \(K\), also for negative \(a\).
\end{proof}

\begin{proof}[Proof of Theorem~\ref{thm:critical}]
Write \(g\) for the prescribed positive nonincreasing envelope
tending to zero, evaluated at \(L(t)=\log(e+|t|)\). Choose open
basis balls \(O_j\) in the torus, repeating each ball infinitely
often, and choose \(\epsilon_j\downarrow0\). In a single dense
forward orbit of \(T=A^{-1}\), find a finite excursion
\(z_{j,0},\ldots,z_{j,\ell_j}\) starting within \(\epsilon_j\) of
zero, visiting \(O_j\) at index \(i_j\), and ending within
\(\epsilon_j/(1+\|T\|)\) of zero. Such excursions exist by density
of all tails. Place the visit at a time \(m_j\in S\) and put
\(N_j=m_j-i_j\). Since \(S\) is unbounded, the windows
\([N_j,N_j+\ell_j]\) can be chosen successively disjoint, with
\[
 N_j\ge j(\ell_j+1),\qquad N_j/j\ge j,\qquad
 g(N_j/j)\le2^{-j}.
\]

Put the excursions at these windows to define a sequence \(z_n\)
on the torus, and put \(z_n=0\) at all other integer times. Write
\(z_{n+1}=Tz_n+\eta_n\pmod{\mathbb Z^2}\) with \(\eta_n\) a
vector lift of the jump error. The errors are bounded, vanish
between jumps, and tend to zero as \(n\to+\infty\).
Use superscripts \(c,e\) for the contracting and expanding
components for \(T\), and let \(T_c,T_e\) be the corresponding
restrictions. Define
\[
 w_n^c=-\sum_{r\ge0}T_c^r\eta_{n-1-r}^c,\qquad
 w_n^e=\sum_{r\ge0}T_e^{-r-1}\eta_{n+r}^e.
\]
The series converge absolutely by geometric bounds and retain the
signs of the eigenvalues. Reindexing gives
\(w_{n+1}=Tw_n-\eta_n\) for \(w_n=w_n^c+w_n^e\).
Geometric domination and \(\eta_n\to0\) give \(w_n\to0\) as
\(n\to+\infty\). Therefore \(q_n=z_n+w_n\pmod{\mathbb Z^2}\)
is the orbit \(T^n y\) of \(y=q_0\). It tends to zero outside
the excursion windows. At the visit times \(m_j\in S\), the
corrections tend to zero. Each nonempty open set contains the
closure of a basis ball visited infinitely often, so \(q_n\) is
dense along \(S\).

Fix any transverse affine segment and any critical widths and
unstable exponent. By \eqref{eq:critical-open-box} and
Proposition~\ref{prop:transverse-slices}, its fiber is dense
\(G_\delta\) of dimension zero. Critical hits infinitely often
outside the windows force
\(p+xv\in[-c_s,c_s]e_s\pmod{\mathbb Z^2}\) by
\eqref{eq:inverse-slice-widths} and \(q_n\to0\) there. This is a
finite set of parameters, since only finitely many lifts of the
stable segment meet the given bounded transverse segment.
If \(|\widehat\mu(t)|=O(g(L(t)))\), then \(\mu\) is atomless.
Averaging \(e^{2\pi i tx}\widehat\mu(t)\) over \([-R,R]\) gives
\(\mu(\{x\})\) in the limit by dominated convergence, and the
decay makes the average tend to zero. For infinitely many hits
inside the windows, Lemma~\ref{lem:parameter-block} applies for
all sufficiently large \(j\), since \(\ell_j/N_j\to0\).
The mass bound is \(O_\mu(e^{-\eta N_j}+g(c_1N_j))\) for some
\(c_1>0\). It is summable because \(c_1N_j\ge N_j/j\) eventually.
Borel--Cantelli proves zero mass for the whole fiber. The choice
of \(y\) preceded the choices of segment, affine origin, and
critical dimensions of the rectangles. Inscribed and circumscribed
critical rectangles give the ball assertions.

For the qualitative assertion, restrict to a nondegenerate closed
interval inside any prescribed open parameter subinterval. The
fiber is comeager there, and that closed interval is the support
of a Rajchman measure, namely interval Lebesgue measure.
Matheron and Zelený's category result implies that the fiber
carries a nonzero Rajchman measure
\cite[Corollary~4.5]{MZ}. Take a compact subset of positive
measure and restrict and normalize. Restriction preserves Rajchman
decay, by approximation in \(L^1\) of the restricting density by
finite linear combinations of characters. This gives a compactly
supported Rajchman probability in each open parameter subinterval.
\end{proof}

\subsection{Stable slices and affine origins}

If \(v=v_s e_s\) with \(v_s\ne0\) and \(p=0\), the iterated segment
converges uniformly to zero. For a fixed nonzero center the fiber
is empty. For \(y=0\) it is \(I\) when \(\beta<\chi\), the set
\(\{x\in I:|xv_s|<c_s\}\) when \(\beta=\chi\), and \(I\cap\{0\}\)
when \(\beta>\chi\). These statements follow from the exact forward
stable displacement \(b^n x v_s e_s\). For balls the exponent
\(\beta\) is the radius exponent and the critical test is
\(|xv_s|<c\).

At a general affine origin, any sufficiently late hit forces
\(y_n-A^n p\) modulo the lattice to have norm less than \(1/4\).
If it has such a short lift \(d_n\), the exact rectangle test at
all sufficiently late times is
\[
 |(d_n)_u|<c_u e^{-\alpha n},\qquad
 \left|x-\frac{b^{-n}(d_n)_s}{v_s}\right|
 <\frac{c_s}{|v_s|}\,e^{-(\beta-\chi)n}.
\]
The sign of \(b^{-n}\) enters the center of this window. At
criticality the window has fixed width, and both its availability
and its location depend on the affine base orbit.

In the target variable, each swept ball neighborhood for a stable
segment lies in a single ball about \(A^n p\) of radius
\(O(\lambda^{-n}+r_n)\), where \(r_n=c\lambda^{-\tau n}\).
Thus \(V\), and hence \(D,P\), have Hausdorff dimension zero.
If \(A^n p\) is dense along \(S\), each open parameter subinterval
gives a dense target tail union, since one fixed parameter in that
subinterval has forward images within \(o(1)\) of \(A^n p\).
Consequently \(D\) is dense \(G_\delta\) in the torus and has local
compact Rajchman probabilities by \cite[Corollary~4.5]{MZ} and
Fourier localization. When \(p=0\) and \(\tau>1\), \(D\) is empty
and \(P=V=\{0\}\) if \(0\in I\), while \(P=V=\varnothing\) if
\(0\notin I\).

\section{A cubic example and the limits of the planar law}
\label{sec:higher}

\begin{example}
The companion matrix of \(X^3-3X^2+1\) has three real eigenvalues, one in each of \((-1,0)\), \((0,1)\), and \((2,3)\), by the signs at the endpoints. The rational root test gives irreducibility. No two roots have the same modulus, since a pair of opposite roots would force the third root to be \(3\) by the sum of the roots. The expanding root is a Pisot unit. The product of the two contracting moduli is the reciprocal of the expanding root, so each contraction exponent is strictly less than \(1\) on the logarithmic scale of the expanding root.

For every \(\tau>1\), consider lines with nonzero expanding component. A line using exactly one contracting eigendirection has the occupied-block criterion from the planar unstable line. A line using both contracting eigendirections has the individual-level criterion from the planar mixed line. In full time these criteria give the respective attainable bounds \(o((\log L(\xi))^{-1})\) and \(o(L(\xi)^{-1})\), where \(L(\xi)=\log(e+|\xi|)\). Increasing either decay exponent from \(1\) to \(1+\varepsilon\), \(\varepsilon>0\), in the corresponding big-\(O\) bound forces zero \(V\)-mass. The inverse companion matrix illustrates the theorem with two expanding eigenvalues at the stated large rates.
\end{example}

\begin{theorem}\label{thm:higher}
Let \(A\in\mathrm{GL}_d(\mathbb Z)\), \(d\ge3\), have an irreducible characteristic polynomial with all roots real, distinct in modulus, and off the unit circle. Choose unit eigenvectors \(e_i\) with eigenvalues \(\lambda_i\), ordered by decreasing modulus. Set \(\Lambda=|\lambda_1|\), \(a_i=\log_{\Lambda}|\lambda_i|\), \(b_*=-a_d\), and \(\sigma=(1-a_2)/(4d)\). Consider the affine segment \(\gamma(x)=z+xv\pmod{\mathbb Z^d}\), \(x\in[0,1]\), with \(v=\sum_i v_i e_i\) and \(v_1\ne0\). Use the same sets \(D\subset P\subset V\) of target centers, defined on \(\mathbb T^d\) with open balls of radius \(c\Lambda^{-\tau n}\), \(c>0\), along any infinite schedule \(S\subseteq\mathbb N\).

With exactly one expanding eigenvalue assume \(\tau>b_*\). With several expanding eigenvalues assume
\[
 \tau>\max\left\{b_*,\frac{b_*}{(d-1)\sigma},\frac1{d-1}\right\}.
\]
All three sets have Hausdorff dimension \(s_0=1+1/\tau\) in every nonempty open target patch and Fourier dimension zero. For a regular gauge \(h\), the criterion is
\[
\begin{cases}
 \displaystyle\sum_{n\in S}h(n), & v_i\ne0\text{ for every }i,\\[4pt]
 \displaystyle\sum_{\substack{j\ge0:\ S\cap[2^j,2^{j+1})\ne\varnothing}} h(2^j), & v_i=0\text{ for some }i.
\end{cases}
\]
If the relevant series converges, every finite positive torus measure with \(|\widehat\mu(k)|=O(h(L(k)))\) on integer frequencies satisfies \(\mu(V)=0\). If it diverges, each nonempty open target patch contains a compact subset of \(D\) supporting a probability with \(|\widehat\mu(\xi)|=o(h(L(\xi)))\) at real frequencies in a local lift and finite \(s\)-energy for every \(0<s<s_0\), simultaneously.
\end{theorem}

The proof is in Appendix~\ref{app:higher}. The rate condition with several expanding directions is sufficient. The optimal range for several expanding directions and the case of an irreducible polynomial with complex conjugate roots remain open here.

\begin{example}[A reducible automorphism]\label{ex:higher-reducible}
Let \(C=\left(\begin{smallmatrix}2&1\\1&1\end{smallmatrix}\right)\) be the cat matrix, and use \(A=\operatorname{diag}(C,C)\) on \(\mathbb T^4\). A through-origin segment with direction \((v',0)\), where \(v'\in\mathbb R^2\) has a nonzero expanding component for \(C\), stays in \(\mathbb T^2\times\{0\}\) at every time. For any radii tending to zero, \(V\) is contained in this rational subtorus. Every nonzero positive measure carried by \(V\) has Fourier coefficient equal to its mass on every character \((0,k')\), \(k'\in\mathbb Z^2\). Consequently there is no probability carried by \(V\) with Fourier transform tending to zero at infinity, including in a local lift.
\end{example}

\subsection{Non-unimodular planar saddles}\label{sec:nonunit}

Let \(A\) be a nonsingular integer \(2\times2\) matrix with signed
eigenvalues \(a,b\) satisfying \(0<\rho=|b|<1<\lambda=|a|\), and let
\(\Delta=|\det A|=\lambda\rho>1\) and \(\chi=\log\lambda\).
Its characteristic polynomial is irreducible, since \(b\) cannot
be a nonzero rational algebraic integer. We use the same unit
eigenvectors, affine lifts \(p+xv\), and notation \(D,P,V\) as
before, with \(\tau>1\) and \(r_n=c\lambda^{-\tau n}\).
Write \(\mathcal B(S)=\{j:S\cap[2^j,2^{j+1})\ne\varnothing\}\).

\begin{proposition}\label{prop:nonunit-gauge}
For a pure unstable line \(v_s=0\), \(v_u\ne0\), the sharp gauge
series is \(\sum_{j\in\mathcal B(S)}h(2^j)\). For a mixed line
\(v_uv_s\ne0\), it is \(\sum_{n\in S}h(n)\). Specifically,
convergence forces \(\mu(V)=0\) for every finite positive torus
measure with \(|\widehat\mu(k)|=O(h(L(k)))\) for a regular gauge
\(h\). Under divergence, every open target patch contains a compact
subset of \(D\) carrying a probability with
\(|\widehat\mu(\xi)|=o(h(L(\xi)))\) at real frequencies in a lift
and finite \(s\)-energy simultaneously for all \(0<s<d_*\), where
\[
 d_*=
 \begin{cases}
 1+1/\tau,&v_s=0,\\[2pt]
 1+(\chi-\tfrac12\log\Delta)/(\tau\chi),&v_uv_s\ne0.
 \end{cases}
\]
In every open patch the Hausdorff dimensions of \(D,P,V\) lie
between \(d_*\) and \(1+1/\tau\), and their Fourier dimensions
are zero.
\end{proposition}

\begin{proof}
The determinant bounds \eqref{eq:norm} hold by irreducibility.
For \(0<|k|\le\lambda^{n/3}\), the unstable term has modulus
\(\gg_v\lambda^{2n/3}\), whereas the stable term is
\(O_v(\rho^n\lambda^{n/3})\). This is the dominance needed
in the Dirichlet strip test of Section~\ref{sec:lines}.
The remaining ingredients there are the bound
\(|A^n v|\ll_v\lambda^n\), and, on the pure line, the single
normal to \(e_u\) on each short block. Both estimates apply
here. Lemma~\ref{lem:majorant} therefore gives the same mixed
level bound \(O_\mu(e^{-\eta n}+h(n))\) and pure block bound
\(O_\mu(e^{-\eta N}+h(N))\), proving convergence.

Use the tube densities \(G_n\) of \eqref{eq:tube} formed with
\(f_J(x)\,dx\) on the segment. The low and high frequency
estimates from \eqref{eq:linefourier} apply by the same dominance
and bump decay. In the middle range \(n\asymp L(k)\), the
pure-line count of selected dyadic representatives is unchanged.
For a mixed line put
\(X=v_u k_u,\ Y=v_s k_s\), and
\(t_*=\log(|Y/X|)/\log(\lambda/\rho)\). The larger term in the
parameter frequency has magnitude
\[
 \sqrt{|XY|}\,\Delta^{n/2}(\lambda/\rho)^{|n-t_*|/2}.
\]
Here \(\sqrt{|XY|}\gg_v1\) and \(\Delta^{n/2}\ge1\). At distance
at least one from \(t_*\), the smaller term is at most
\(\rho/\lambda\) times the larger. Thus the sum of the magnitudes
of the parameter Fourier factors over all \(n\) is bounded
uniformly in \(k\ne0\). Doubling compares \(h(n)\) and
\(h(L(k))\) in the middle range. Thus \eqref{eq:weighted}
holds with weights \(w_n=h(n)\) on all of \(S\) for mixed
lines, and on one representative of each occupied dyadic
block for pure lines.

We give the ball estimates needed for the energy assertion.
For the unsmoothed pushforward, each passage through a ball of
radius \(R\) corresponds to a lattice lift of its center within
\(R\) of the lifted segment. If \(x,x'\) witness two distinct
passages and \(t=|x-x'|\), the product bound yields
\begin{equation}\label{eq:nonunit-passage}
 c_0\le |(z-z')_u(z-z')_s|
 \ll_v \Delta^n t^2+\lambda^nRt+R^2
\end{equation}
for the two center lifts \(z,z'\). For \(R\) sufficiently small
this forces \(t\gg(\Delta^{n/2}+\lambda^nR)^{-1}\). A passage
occupies a parameter interval of length \(O(R\lambda^{-n})\),
by the unstable projection. With bounded parameter density the
ball mass is therefore \(O_J(R^2+R\lambda^{-n}\Delta^{n/2})\).
For a pure line the quadratic term vanishes and the ball mass
is \(O_J(R^2+R\lambda^{-n})\), allowing one passage when
\(\lambda^nR<1\). These bounds extend to \(R\le1\) by increasing
constants.

For \(R\ge r_n\), the same bounds for \(G_n\) follow by
enlarging the ball by \(r_n/4\). In either case the mass is
\(O_J(R^{d_*})\) since the second term has the form
\(R\) times \(O(r_n^{d_*-1})\). At \(R<r_n\) use the density
bound \(O_J(r_n^{d_*-2})\), obtained from the unsmoothed mass
at scale \(r_n\) and the maximum of the bump density. This gives
the same \(O_J(R^{d_*})\) estimate. Apply Lemma~\ref{lem:average}
to the weighted coefficient bound and these uniform ball
bounds. Divergence gives arbitrarily large total weight in
every tail. Lemma~\ref{lem:compact} then gives the probability
on \(D\), with all the claimed energies, by cycling the parameter
intervals \(J\).

The \(r_n\)-neighborhood of the segment is covered by
\(O(\lambda^n/r_n)\) balls of radius \(2r_n\), proving the
dimension upper bound. For each infinite \(S\) one can choose a
regular gauge with divergent occupied-block sum. Group successive
dyadic indices into finite consecutive ranges with at least \(2^j\)
occupied indices in the \(j\)-th range, and put
\(h(t)=2^{-j}\) when the dyadic index of \(t\) is in that range.
Doubling crosses at most one boundary, so this is a regular
gauge. Both series then diverge and the energy lower bound
gives the local Hausdorff lower bound for every \(S\).
Every positive power of Fourier decay implies the convergent
bound for \(h(t)=t^{-2}\), giving Fourier dimension zero.
\end{proof}

\begin{corollary}\label{cor:nonunit-fiber-bound}
For the same supercritical balls, every fixed-center transverse
fiber on a mixed line has Hausdorff dimension at most
\(\log\Delta/(2(\tau+1)\chi)\). On a pure unstable line it has
dimension zero.
\end{corollary}

\begin{proof}
With \(R=r_n\), \eqref{eq:nonunit-passage} gives
\(O(\Delta^{n/2})\) passages, each a window of length
\(O(r_n\lambda^{-n})\). The pure-line count is bounded.
The Hausdorff sum gives the stated upper bounds.
\end{proof}

\begin{proposition}\label{prop:rational-winding}
If \(v=t w\), where \(w\) is a primitive integer vector and
\(t\ne0\), put \(g_n=\gcd((A^nw)_1,(A^nw)_2)>0\) and
\[
 \gamma_A=\sum_{\mathfrak p\mid\Delta}
 \min\left\{v_{\mathfrak p}(\operatorname{tr} A),
            \tfrac12 v_{\mathfrak p}(\det A)\right\}
 \log\mathfrak p ,
\]
where the sum is over primes and \(v_{\mathfrak p}(0)=+\infty\).
Then \(\log g_n=n\gamma_A+O_{A,w}(1)\). In every open target
patch the Hausdorff dimensions of \(D,P,V\) all equal
\begin{equation}\label{eq:rational-winding-dimension}
 d_{\mathrm{rat}}=1+\frac{\chi-\gamma_A}{\tau\chi}.
\end{equation}
For the divergent mixed series the probability in
Proposition~\ref{prop:nonunit-gauge} can be taken to have finite
\(s\)-energy for all \(0<s<d_{\mathrm{rat}}\).
\end{proposition}

\begin{proof}
Cyclicity of \(w\) over \(\mathbb Q\) gives the comparison
\(\|A^n w\|_{\mathfrak p}\asymp_{A,w,\mathfrak p}
\|A^n\|_{\mathfrak p}\) in entrywise maximum norms.
Indeed, \(B_w=[w,Aw]\) is invertible and
\(A^n B_w=[A^nw,A(A^nw)]\).
Diagonalize over a finite \(\mathfrak p\)-adic splitting field.
The two roots are distinct, so the matrix norm is comparable
to \(\max(|a|_{\mathfrak p},|b|_{\mathfrak p})^n\).
This gives
\[
 v_{\mathfrak p}(g_n)
 =n\min\{v_{\mathfrak p}(a),v_{\mathfrak p}(b)\}+O(1).
\]
For unequal valuations the smaller is the trace valuation.
For equal valuations their common value is half the determinant
valuation and the trace valuation is at least as large.
Only primes dividing \(\Delta\) can divide \(g_n\), since \(w\)
is primitive and \(A\) is invertible modulo every other prime.
Summation gives the formula for \(\log g_n\). For the gcd of the
entries of matrix powers, Noferini identifies the slope with the
minimum eigenvalue valuation and proves eventual periodicity of the
valuation increments \cite[Theorems~2.4--2.5]{Noferini}.

The primitive vector \(w_n=A^nw/g_n\) has length comparable
to \(e^{(\chi-\gamma_A)n}\), since \(w_u\ne0\). In particular
\(0\le\gamma_A\le\tfrac12\log\Delta<\chi\).
The closed geodesic in direction \(w_n\) has lifts spaced by
\(1/|w_n|\), from the primitive integer normal to \(w_n\).
Its normalized length measure gives a ball of radius \(R\)
mass \(O(R^2+R/|w_n|)\). Indeed, at most
\(O(1+R|w_n|)\) lifts meet the ball and each chord has length
at most \(2R\).
The pushforward of a fixed \(f_J\) is dominated by a constant
times normalized length measure. The circle parameter is
\(t g_n x\bmod 1\), whose density is bounded uniformly because
\(|t|g_n\ge |t|>0\). The smoothing argument in
Proposition~\ref{prop:nonunit-gauge} now gives the ball bound
\(O_J(R^{d_{\mathrm{rat}}})\), and the same Fourier estimates
and construction give the asserted energies. The tube of the
whole geodesic is covered by \(O(|w_n|/r_n)\) balls of radius
\(2r_n\), giving the upper bound. This proof works for any affine
origin, nondegenerate segment, and infinite schedule.
\end{proof}

For the whole closed geodesic through the origin the earlier
restricted-linear-form results also give the dimensions of \(V\).
If \(q_n\) is the primitive integer normal, its width in the
linear-form test is \(|q_n|r_n\). Squeezing this width by powers
with exponents \(1-\tau\chi/(\chi-\gamma_A)\pm\varepsilon\)
and using convergence exponent zero for the exponentially
growing \(|q_n|\) gives
\eqref{eq:rational-winding-dimension} by Rynne's theorem
\cite[p.~70]{Rynne}. The upper-bound proof of Cai and Hambrook
\cite[Theorem~1.4.1 and \S3.3]{CH} gives Fourier dimension zero, since
the widths sum and \(\sum_{n\in S}r_n^s<\infty\) for all \(s>0\).
If \(\gcd(\operatorname{tr}A,\Delta)=1\), then \(\gamma_A=0\)
and every rational direction has the full dimension \(1+1/\tau\).
For general irrational mixed directions the exact Hausdorff
dimension remains open here.

\begin{example}\label{ex:nonunit-cat}
Take \(A=2A_{\mathrm{cat}}\) where
\[
 A_{\mathrm{cat}}=\begin{pmatrix}2&1\\1&1\end{pmatrix},
 \qquad \lambda_{\mathrm{cat}}=(3+\sqrt5)/2 .
\]
Here \(\lambda=2\lambda_{\mathrm{cat}}\), \(\rho=2/\lambda_{\mathrm{cat}}<1\),
and \(\Delta=4\). For \(w=(1,0)\), \(g_n=2^n\) exactly,
since \(A_{\mathrm{cat}}^nw\) stays primitive.
The rational dimension in \eqref{eq:rational-winding-dimension}
attains the mixed lower bound in Proposition~\ref{prop:nonunit-gauge}
and is strictly smaller than \(1+1/\tau\).

For \(p=y=0\) and \(v=w\), the bound in
Corollary~\ref{cor:nonunit-fiber-bound} is attained too.
Writing \(w_n=A_{\mathrm{cat}}^n w\), every lattice point off
\(\mathbb R w_n\) is at distance at least \(1/|w_n|\) from
that line. As \(r_n|w_n|\to0\), the hits for all large \(n\)
are exactly
\[
 \|2^n x\|_{\mathbb R/\mathbb Z}<r_n/|w_n|.
\]
The parameter intervals have spacing \(2^{-n}\) and radii
comparable to \(\lambda^{-(\tau+1)n}\). They give
\[
 \dim_H E_0=\frac{\log 2}{(\tau+1)\log\lambda}
\]
for every infinite \(S\). The upper bound is the Hausdorff sum.
For the lower bound, raising the radii to this power gives a
fixed positive multiple of the spacing. The enlarged intervals
then cover a fixed positive fraction of every subinterval at
each sufficiently late \(n\in S\). The density theorem gives
full measure for their limsup, and mass transference gives the
lower bound \cite[Theorem~2]{BV}.

The stable multiplier gives the contraction threshold \(r_n=c\rho^n\).
For this map it occurs at \(\tau=-\log\rho/\chi<1\).
Fix any center \(y\) and take
\(p=0,\ v=w\). The distance from \(y\) to the closed geodesic
in direction \(w_n\) is \(O(|w_n|^{-1})=o(r_n)\). Let \(u_n\)
be the circle coordinate of a nearest point. For a fixed
sufficiently small \(C>0\), hits contain all the windows
\(\|2^n x-u_n\|_{\mathbb R/\mathbb Z}<C r_n/|w_n|\) at large
times. For every regular gauge \(g\) with
\(\sum_{j\in\mathcal B(S)} g(2^j)=\infty\), these windows give a
compactly supported probability in the fiber in every open
parameter subinterval, with real-frequency decay \(o(g(L))\).
In particular, on full times a single such gauge works for every
\(y\), so the prescribed-envelope obstruction of
Theorem~\ref{thm:critical} fails for this map at its contraction
threshold.

To prove the decay claim, take a smooth unit-integral bump on
the circle supported within distance
\(\delta_n=C r_n/(2|w_n|)\) of zero, obtained by rescaling one
fixed profile. Evaluate it at \(2^n x-u_n\) and denote the
result by \(P_n(x)\). Its constant coefficient is one, its
other frequencies are nonzero multiples of \(2^n\), and the
absolute coefficient sum is \(O(\delta_n^{-1})\). For a fixed
smooth compactly supported probability density \(f\) in the
parameter interval, this gap and rapid decay of \(\widehat f\)
give \(\widehat{fP_n-f}(s)=O_{f,M}(e^{-Mn})\) for every \(M\)
on \(|s|\le e^{\kappa n}\), \(0<\kappa<\log2\). At
\(|s|\ge e^{Kn}\) for \(K\) sufficiently large the difference
is \(O_f((1+|s|)^{-2})\) by repeated integration by parts.
Indeed the order-\(j\) derivative is bounded in \(L^1\)
by \(O_{f,j}(\delta_n^{-1}(2^n/\delta_n)^j)\).
In between, positivity and the mass estimate at zero give a
constant bound.

Select one time from each occupied dyadic block with weights
\(g(n)\). Only boundedly many selected times satisfy
\(e^{\kappa n}<|s|<e^{Kn}\), and there \(g(n)\asymp g(L(s))\).
The two tails are summable with error \(O_f(g(L(s)))\), by the
exponential estimates and doubling. Therefore normalized
multiplication of \(f\) by a finite weighted average of the
\(P_n\) changes its Fourier transform in the weighted supremum
norm by \(O_f(W^{-1})\), where \(W\) is the total weight.
Normalization preserves this estimate by evaluating first at
zero. Divergence allows arbitrarily large \(W\) in each tail.
Iterating with summable errors gives nested compact supports
in successive finite unions of hit windows at advancing times.
The weak limit has support in the fiber. The weighted Fourier
transforms converge uniformly, and the quotient for every
smooth stage vanishes at infinity. This proves the little-\(o\)
bound for the limit.
\end{example}

\subsection{Questions at the boundary}
The sparse straight-line criteria classify every regular gauge. For a
curved arc an arbitrary schedule may lose divergence of
\(\sum h(n)^{2/3}\) on passage to logarithmically separated times.
It remains to decide whether divergence of the unthinned series suffices.
Even at full time the estimates leave open the refinements
\(h(t)=t^{-3/2}(\log(e+t))^{-q}\), \(0<q\le3/2\).

The critical center in Theorem~\ref{thm:critical} is chosen for a
specified envelope. The optimal parameter-space decay for a Haar-typical
critical center remains undetermined. For a non-unimodular saddle the
dimension of the center set for an arbitrary irrational mixed direction
is another open case of these estimates. Finally a nonreal conjugate pair
introduces rotating phases of equal modulus. The distinct-real-modulus
argument below leaves the sharp mixed-direction gauge for such spectra open.

\appendix
\section{Proof of Theorem~\ref{thm:higher}}
\label{app:higher}

Write \(\rho_n=\Lambda^{-\tau n}\), so the target radius is \(c\rho_n\). Throughout, constants may depend on the fixed map, segment, and radius parameters. We use the right eigenbasis for positions and write \(k_i=k\cdot e_i\) for the dual pairings. In particular, \(k\mapsto(k_i)_{i=1}^d\) is a fixed invertible real linear change of coordinates.

\subsection{Algebraic norm and the frequency gap}

Choose a nonzero right eigenvector for \(\lambda_1\) with entries in \(K=\mathbb Q(\lambda_1)\), and multiply by an integer so its entries are algebraic integers. Each embedding of \(K\) gives a real conjugate eigenvector, and each eigendirection occurs once since the characteristic polynomial is irreducible. The resulting columns are independent. If their pairing with an integer \(k\) vanished for one embedding, the pairing would vanish for every embedding, giving \(k=0\). For \(k\ne0\) the product of the unnormalized pairings is thus a nonzero integer, the field norm of an algebraic integer. Normalizing the columns yields
\begin{equation}\label{eq:higher-norm}
 \prod_{i=1}^d |k_i|\gg 1,\qquad
 |k_1|\gg |k|^{-(d-1)}\quad(0\ne k\in\mathbb Z^d).
\end{equation}
Since \(\sum_i a_i=0\), we have \(a_2\ge-1/(d-1)\). Consequently \(1-2\sigma(d-1)>0\), and \(1-2\sigma(d-1)-(a_2+2\sigma)=(1-a_2)/2>0\). For \(0<|k|\le\Lambda^{2\sigma n}\), the first term of \(k\cdot A^n v\) has modulus \(\gg\Lambda^{[1-2\sigma(d-1)]n}\) by \eqref{eq:higher-norm} and \(v_1\ne0\). The rest have total modulus \(O(\Lambda^{(a_2+2\sigma)n})\). Hence for all sufficiently large \(n\) there is a fixed \(\kappa>0\) such that
\begin{equation}\label{eq:higher-gap}
 |k\cdot A^n v|\gg\Lambda^{\kappa n}
 \qquad (0<|k|\le\Lambda^{2\sigma n}).
\end{equation}
This comparison uses moduli and holds with either sign of each eigenvalue.

\subsection{Upper bounds for target mass}

In both rate ranges \(\tau>1/(d-1)\). In the one-expanding case this follows because the \(d-1\) contraction exponents sum to \(1\). Fix \(1/(d-1)<t<\tau\) and \(0<\eta<(d-1)t-1\). In a real orthonormal frame with first axis parallel to \(A^n v\), use the centered box with first half-width \(\Lambda^{-\eta n}/|A^n v|\) and the other half-widths \(\Lambda^{tn}\). Since \(|A^n v|\asymp\Lambda^n\), its volume grows as \(\Lambda^{((d-1)t-1-\eta)n}\) up to constants. Minkowski's convex body theorem gives a nonzero integer point \(k\) for large \(n\), with \(|k|\ll\Lambda^{tn}\) and \(|k\cdot A^n v|\le\Lambda^{-\eta n}\). By \eqref{eq:higher-gap}, \(|k|>\Lambda^{2\sigma n}\). Each \(y\in U_n([0,1])\) satisfies
\begin{equation}\label{eq:higher-one-phase}
 \|k\cdot y-k\cdot A^n z\|_{\mathbb R/\mathbb Z}
 \le \Lambda^{-\eta n}+c\rho_n|k|\ll e^{-\beta n}
\end{equation}
for some \(\beta>0\). The single-phase instance of Lemma~\ref{lem:majorant} gives mass \(O(e^{-\beta' n}\|\mu\|+\sup_{\ell\ne0}|\widehat\mu(\ell k)|)\) for some \(\beta'>0\). Under the stated Fourier bound this is \(O(e^{-\beta' n}\|\mu\|+h(n))\), using the lower bound on \(|k|\), eventual monotonicity and doubling. The sum over \(S\) gives the individual-level converse. The same test applies to any direction under consideration, so a positive power bound on \(|\widehat\mu(k)|\) makes the masses summable for every schedule and proves zero Fourier dimension.

For the occupied-block converse, put \(I_v=\{i:v_i\ne0\}\) and suppose \(I_v\ne\{1,\ldots,d\}\). Fix \(b_*<t<\tau\). Choose \(\eta>0\) small enough that \(\sum_{i\notin I_v}(a_i+t)-|I_v|\eta>0\). In the dual coordinates take the box
\[
 |k_i|\le
 \begin{cases}
 \Lambda^{(-a_i-\eta)N},&i\in I_v,\\
 \Lambda^{tN},&i\notin I_v .
 \end{cases}
\]
The exponent of its volume is the positive number just displayed, by \(\sum_i a_i=0\). The coordinate change contributes only a fixed determinant. Minkowski gives a nonzero integer vector for large \(N\). Since every \(-a_i\le b_*<t\), its norm is \(O(\Lambda^{tN})\). Also \(1\in I_v\), so \eqref{eq:higher-norm} and \(|k_1|\le\Lambda^{-(1+\eta)N}\) give \(|k|\gg\Lambda^{(1+\eta)N/(d-1)}\).

Choose \(\delta>0\) with \(\delta\max_i |a_i|\le\eta/2\). For every integer \(N\le n\le(1+\delta)N\) the cap on each active coordinate gives \(|k_i\lambda_i^n|\le\Lambda^{-\eta N/2}\), and \(c\rho_n|k|\ll\Lambda^{-(\tau-t)N}\). Thus the phase test in \eqref{eq:higher-one-phase} holds throughout the block with error \(\epsilon_N\ll e^{-\beta N}\), where \(\beta>0\) is independent of further decreases in \(\delta\). The phases \(\theta_n=k\cdot A^nz\bmod 1\) can vary with \(n\).

To apply Lemma~\ref{lem:majorant}, let \(M=\lfloor e^{\beta N/2}\rfloor\) and let \(m\le\delta N+1\) be the number of times tested in the block. The vectors \((q\theta_n)_n\) modulo one, \(1\le q\le M\), fall into \(16^m\) boxes of side \(1/16\). In the most populated box the indices form a set \(F\) of size at least \(M/16^m\). Taking \(\delta\log16<\beta/4\) makes \(|F|\) exponentially large. On any of the tests \(\|u-\theta_n\|_{\mathbb R/\mathbb Z}\le\epsilon_N\), the points \(qu\bmod1\) for \(q\in F\) lie within \(M\epsilon_N\) of an arc of length \(1/16\). Thus \(\big|\,|F|^{-1}\sum_{q\in F}e^{2\pi iqu}\big|^2\) has a fixed positive lower bound there. This polynomial has constant coefficient \(|F|^{-1}\) and absolute coefficient sum \(1\). Applied to \(u=k\cdot y\) and integrated, it bounds the mass of the union by
\[
 O\!\left(e^{-\beta' N}\|\mu\|+
             \sup_{\ell\ne0}|\widehat\mu(\ell k)|\right)
 =O\!\left(e^{-\beta' N}\|\mu\|+h(N)\right)
\]
for some \(\beta'>0\). Each dyadic block \([2^j,2^{j+1})\) is a union of \(O_\delta(1)\) such short blocks or their truncations, with \(N\asymp2^j\). Keeping only the occupied dyadic blocks gives a summable bound when the block series converges. In both converses the tail-sum estimate gives zero mass to the limsup \(V\).

\subsection{Weighted level densities}

Fix an open parameter interval \(J\) from a countable basis in \((0,1)\). Choose a smooth nonnegative probability density \(f_J\) on \(\mathbb R\), compactly supported in \(J\), and a smooth nonnegative probability bump \(\chi\) on \(\mathbb R^d\) supported in \(B(0,c/4)\). Let \(G_{n,J}\) be the density obtained by pushing \(f_J(x)\,dx\) forward by \(x\mapsto A^n(z+xv)\bmod\mathbb Z^d\) and convolving with the periodization of \(\rho_n^{-d}\chi(\cdot/\rho_n)\). It has integral one and compact support in \(U_n(J)\). Its Fourier coefficients are
\begin{equation}\label{eq:higher-density}
 \widehat G_{n,J}(k)
 =e^{-2\pi i k\cdot A^nz}\widehat\chi(\rho_n k)
                   \widehat f_J(k\cdot A^n v).
\end{equation}
For \(0<|k|\le\Lambda^{2\sigma n}\), the gap \eqref{eq:higher-gap} and rapid decay of \(\widehat f_J\) give \(O_B(\Lambda^{-Bn})\) for any \(B>0\). For \(|k|\ge\Lambda^{(\tau+1)n}\) we have \(\rho_n|k|\ge |k|^{1/(\tau+1)}\), so the bump factor gives \(O_B(|k|^{-B})\). For a fixed \(k\ne0\) there are \(O(L(k))\) times in the latter range, and the low-frequency estimate sums over a geometric tail starting at \(n\gg L(k)\) for large \(|k|\). A regular gauge has a polynomial lower bound as a function of its argument, by doubling. With bounded weights the sum of the two tails is therefore \(O(h(L(k)))\).

In the intermediate range \(n\asymp L(k)\), doubling gives \(h(n)\asymp h(L(k))\). If a component of \(v\) vanishes, choose one representative from \(S\) in each occupied dyadic block. Only boundedly many representatives can lie in the intermediate range for a fixed \(k\). Give each representative weight \(h(n)\). The coefficient bound \(|\widehat G_{n,J}(k)|\le1\) handles this range.

If all components of \(v\) are nonzero, use every \(n\in S\) with weight \(h(n)\). For fixed \(k\ne0\) put \(T_i(n)=|v_i k_i|\Lambda^{a_i n}\) and \(M(n)=\max_i T_i(n)\). The product of the \(T_i(n)\) is uniformly bounded below by a positive constant by \eqref{eq:higher-norm} and \(\sum_i a_i=0\). For each pair \(i\ne j\), the condition \(T_i(n)/T_j(n)\in[1/(2d),2d]\) holds on a real interval of bounded length, since the slope of the logarithm of that ratio is nonzero and independent of \(k\). Thus outside \(O_A(1)\) integer times the largest term exceeds each other by a factor \(2d\), giving
\[
 |k\cdot A^n v|\ge M(n)/2 .
\]
If \(M(n)\le H\), \(H\ge2\), then \(T_1(n)\in[c_v H^{-(d-1)},H]\) for a fixed \(c_v>0\). As \(T_1(n)\) grows by a factor \(\Lambda\) each time, this accounts for \(O_{A,v}(1+\log H)\) integer times. Summing the rapid decay of \(\widehat f_J\) on dyadic ranges of \(M(n)\), and bounding the exceptional times by \(1\), yields \(\sum_{n\ge1}|\widehat f_J(k\cdot A^n v)|\ll_{A,J,v}1\) uniformly in \(k\ne0\). This controls the weighted intermediate range in the all-component case as well.

Writing \(S'\) for the selected times in either case, we have obtained
\begin{equation}\label{eq:higher-average}
 \sum_{n\in S'}h(n)|\widehat G_{n,J}(k)|
       \ll_{A,J,v,h} h(L(k))\qquad(k\ne0).
\end{equation}
The zero mode is exactly \(1\) at each level, as required for the normalized averages of Lemma~\ref{lem:average}. In the block case the sum of the weights on \(S'\) diverges precisely when the occupied-block series does, by doubling.

\subsection{Volume bounds and the compact construction}

Let \(\nu_{n,J}\) be the pushforward of \(f_J(x)\,dx\) before smoothing. We first prove \(\nu_{n,J}(B(y,R))\ll_J R^{s_0}\) for \(\rho_n\le R\le1\). In the lift, consider the lattice translates of an \(R\)-ball that meet the swept segment. Their translation vectors lie in an eigenbasis box of widths \(O(|\lambda_i|^n+R)\). The cells obtained by adding \(A^n[0,1)^d\) to these lattice vectors have disjoint interiors because \(A\) is unimodular. Each has volume one and eigenbasis extents \(O(|\lambda_i|^n)\). The enlarged box thus has volume \(O(\prod_i(|\lambda_i|^n+R))\), bounding the number of translates. Each translate contributes \(O_J(R\Lambda^{-n})\) in parameter mass since \(|A^n v|\asymp\Lambda^n\).

With one expanding eigenvalue, set \(b_i=-a_i>0\) for \(i>1\). The product bound just proved gives mass
\[
 \ll_J R\Lambda^{-n}\prod_{i>1}(1+R\Lambda^{b_i n}).
\]
For \(R=\Lambda^{-tn}\), \(0\le t\le\tau\), this is \(O_J(R\Lambda^{-n\sum_{i>1}\min(b_i,t)})\), because \(\sum_{i>1}b_i=1\). Since \(b_i\le b_*<\tau\), each \(\min(b_i,t)\ge b_it/\tau\). The resulting bound is \(O_J(R^{s_0})\).

With several expanding eigenvalues, write \(R_0=\Lambda^{-\sigma n}\). At scales \(R_0\le R\le1\), we show an \(O_J(R^d)\) estimate, also for the pushforward of a smooth nonnegative compactly supported density \(g\) on \(\mathbb R\) majorizing the indicator of \([-1,2]\). Use a periodized smooth ball majorant with mean \(O(R^d)\) and coefficients \(O_B(R^d(1+R|k|)^{-B})\). Such a majorant follows by translating and dilating a fixed bump which is at least one on the unit ball. Up to \(\Lambda^{2\sigma n}\), the nonzero coefficients of the pushforward measure are \(O_M(\Lambda^{-Mn})\) for every \(M>0\), by \eqref{eq:higher-gap} and smoothness. The number of these frequencies is \(O(\Lambda^{2d\sigma n})\). Above \(\Lambda^{2\sigma n}\), use the boundedness of the pushforward coefficients and sum the majorant coefficients, using
\[
 \sum_{|k|>\Lambda^{2\sigma n}}(1+R|k|)^{-B}
 \ll \Lambda^{2d\sigma n}(R\Lambda^{2\sigma n})^{-B}
 \le \Lambda^{(2d-B)\sigma n}
\]
for \(B>d\). Choosing \(M\) and \(B\) sufficiently large proves the asserted ball estimate for both parameter densities. Since \(s_0<d\), it also proves the needed \(R^{s_0}\) bound at these scales.

If \(R<R_0\), each lifted \(R\)-ball that meets the segment with \(x\in[0,1]\) contributes \(\gg R_0\Lambda^{-n}\) of parameter mass for \(g\) in the concentric \(2R_0\)-ball, using an interval around a meeting parameter contained in \([-1,2]\) for large \(n\). Those enlarged balls over distinct lattice translates are disjoint for large \(n\). Their total \(g\)-mass is bounded by the large-scale estimate \(O(R_0^d)\). Thus \(O(\Lambda^nR_0^{d-1})\) original translates can meet the segment, and their total mass for \(f_J\) is \(O_J(RR_0^{d-1})\). When \(R\ge\Lambda^{-b_*n}\), the rate assumption gives \(R_0^{d-1}\le R^{1/\tau}\). When \(R\le\Lambda^{-b_*n}\), every \(|\lambda_i|^n\ge R\), so the cell-volume bound gives \(O(1)\) translates by \(\prod_i|\lambda_i|^n=1\). Their mass \(O_J(R\Lambda^{-n})\) is \(O_J(R^{s_0})\) for \(R\ge\rho_n\). Enlarging constants handles finitely many initial levels.

For the smoothed density \(G_{n,J}\), ball mass at \(R\ge\rho_n\) is bounded by the unsmoothed mass at \(R+c\rho_n/4\), hence is \(O_J(R^{s_0})\). At smaller \(R\), the unsmoothed estimate at a fixed multiple of \(\rho_n\) and the bump's sup norm give density \(O_J(\rho_n^{s_0-d})\). Ball mass is then \(O_J(R^d\rho_n^{s_0-d})\), which is \(O_J(R^{s_0})\). All these estimates are uniform in the center \(y\).

Now take a regular gauge satisfying the corresponding divergence condition. Each finite weighted mean \(Q\) from a fresh tail of \(S'\) has the same ball bounds. Its mass is one, and \eqref{eq:higher-average} bounds the weighted Fourier norm of \(Q-1\) by a constant divided by the total weight. For \(0<s<s'<s_0\), integrating the ball masses against \(R^{-s'-1}\,dR\) gives a uniform bound on the \(s'\)-energy. In the torus energy norm of Lemma~\ref{lem:average} this bounds the contribution of \(|k|>K\) to the squared \(s\)-norm of \(Q-1\) by \(O(K^{s-s'})\), using \((1+|k|)^{s-d}\le (1+K)^{s-s'}(1+|k|)^{s'-d}\) there. The low modes tend to zero as the total weight grows. This verifies the averaging hypotheses in dimension \(d\).

Lemma~\ref{lem:compact} now gives a probability in any lifted target chart inside the prescribed open patch. The compact support lies in \(D\), since each support point has hits in every basis interval in every time tail and the corresponding parameter tail unions are open. The same probability has finite energy at each \(s<s_0\) and little-\(o\) decay at real frequencies by the energy iteration and chart localization in that lemma.

Finally the \(n\)-th tube can be covered by \(O(\Lambda^n/\rho_n)\) balls of radius \(O(\rho_n)\) because the lifted segment has length \(O(\Lambda^n)\). The resulting Hausdorff sum converges at every order above \(s_0\). For the lower bound along any infinite schedule, there is a regular gauge for which even the occupied-block sum diverges. Choose successive cutoffs \(2^{j_m}\), \(j_{m+1}\ge j_m+2\), so the intervening range contains at least \(2^m\) occupied dyadic blocks. Assign \(h(t)=2^{-m}\) on that range and a fixed positive value below the first cutoff. Doubling holds since doubling \(t\) crosses at most one cutoff, and each range contributes at least one to the block series. This also implies divergence of the individual-level sum. The compact probability just constructed in any patch gives the dimension lower bound by its energies. The inclusions \(D\subset P\subset V\) complete the dimension claims.

\begin{thebibliography}{99}
\bibitem{AB} D. Allen and S. Baker, \href{https://doi.org/10.1007/s00029-019-0484-9}{A general mass transference principle}, Selecta Math. (N.S.) \textbf{25} (2019).

\bibitem{BV} V. Beresnevich and S. Velani, \href{https://doi.org/10.4007/annals.2006.164.971}{A mass transference principle and the Duffin--Schaeffer conjecture for Hausdorff measures}, Ann. of Math. (2) \textbf{164} (2006), 971--992.

\bibitem{Bluhm} C. Bluhm, \href{https://doi.org/10.1090/S0002-9939-00-05276-X}{Liouville numbers, Rajchman measures, and small Cantor sets}, Proc. Amer. Math. Soc. \textbf{128} (2000), 2637--2640.

\bibitem{CH} T. Cai and K. Hambrook, \href{https://arxiv.org/abs/2403.19410}{On the exact Fourier dimension of sets of well-approximable matrices}, arXiv:2403.19410.

\bibitem{CZ} S. Chow and Q.-L. Zhou, \href{https://arxiv.org/html/2511.14954v1}{Twisted Diophantine approximation for matrix transformations of tori}, arXiv:2511.14954.

\bibitem{FN} R. Fraser and T. Nguyen, \href{https://arxiv.org/abs/2409.02854}{Sharp Fourier decay estimates for measures supported on the well-approximable numbers}, arXiv:2409.02854.

\bibitem{HV} R. Hill and S. L. Velani, \href{https://doi.org/10.1112/S0024610799007681}{The shrinking target problem for matrix transformations of tori}, J. Lond. Math. Soc. (2) \textbf{60} (1999), 381--398.

\bibitem{HPWZ} Z.-N. Hu, T. Persson, W. Wu and Y. Zhang, \href{https://doi.org/10.1112/jlms.70287}{On shrinking targets for linear expanding and hyperbolic toral endomorphisms}. Numbering follows \href{https://arxiv.org/html/2405.02582v1}{arXiv:2405.02582v1}.

\bibitem{Kahane} J.-P. Kahane, \href{https://doi.org/10.1007/978-3-662-59158-1}{S\'eries de Fourier absolument convergentes}, Ergebnisse der Mathematik \textbf{50}, Springer, 1970.

\bibitem{Kaufman} R. Kaufman, \href{https://doi.org/10.4064/aa-39-3-265-267}{On the theorem of Jarn\'ik and Besicovitch}, Acta Arith. \textbf{39} (1981), 265--267.

\bibitem{Lyons} R. Lyons, \href{https://gdz.sub.uni-goettingen.de/download/pdf/PPN356556735_0083/LOG_0050.pdf}{The measure of non-normal sets}, Invent. Math. \textbf{83} (1986), 605--616.

\bibitem{MZ} \'E. Matheron and M. Zelen\'y, \href{https://doi.org/10.4064/fm185-2-1}{Rudin-like sets and hereditary families of compact sets}, Fund. Math. \textbf{185} (2005), 97--116.

\bibitem{Noferini} V. Noferini, \href{https://arxiv.org/html/2511.22814v2}{Eventual periodicity of the Smith forms of integer matrix powers}, arXiv:2511.22814v2.

\bibitem{PR} I. Polasek and E. Rela, \href{https://arxiv.org/abs/2602.15267}{Slowly decaying Rajchman measures and a restriction theorem for the Fourier transform at the limit case of zero Fourier dimension}, arXiv:2602.15267.

\bibitem{PVZZ} A. D. Pollington, S. Velani, A. Zafeiropoulos and E. Zorin, \href{https://doi.org/10.1093/imrn/rnaa307}{Inhomogeneous Diophantine approximation on \(M_0\)-sets with restricted denominators}, Int. Math. Res. Not. IMRN (2022), 8571--8643. Numbered results refer to \href{https://arxiv.org/abs/1906.01151}{arXiv:1906.01151}.

\bibitem{Rynne} B. P. Rynne, \href{https://matwbn.icm.edu.pl/ksiazki/aa/aa61/aa6115.pdf}{The Hausdorff dimension of certain sets arising from Diophantine approximation by restricted sequences of integer vectors}, Acta Arith. \textbf{61} (1992), 69--81.

\bibitem{Stein} E. M. Stein, \href{https://doi.org/10.1515/9781400883929}{Harmonic Analysis: Real-Variable Methods, Orthogonality, and Oscillatory Integrals}, Princeton Math. Ser. \textbf{43}, 1993.

\bibitem{VdC} J. G. van der Corput, \href{https://www.numdam.org/item/CM_1935__1__15_0/}{Zur Methode der station\"aren Phase. Erste Mitteilung: einfache Integrale}, Compositio Math. \textbf{1} (1935), 15--38.

\bibitem{Power} \emph{Sharp Fourier decay for simultaneous power targets and recurrence}, companion manuscript.
\end{thebibliography}
\end{document}
